% !TEX program = xelatex
%% ko/main.tex : 국문 원고(영문 Scientific Reports 원고 paper/manuscript/en/main.tex 의 쌍둥이판), 2026-10-05 조립.
%% 컴파일: ./build.sh (xelatex, bibtex sn-nature.bst, 3회; 로그 build/). 영문판과 같은 그림 PDF, 같은 Table 1 본체,
%% 같은 참고문헌(../en/references.bib)을 쓴다. 프리앰블의 글꼴과 제목 서식은 대회 보고서 outputs/report/main.tex 6-33행에서
%% 가져와 1단으로 바꿨다. \xeCJKsetup{CJKspace=true} 는 한글 어절 사이 공백을 지킨다(scirep_format_and_drafts.md 5.2 시험 5).
%% 순서(MANUSCRIPT_SPEC 1.1, 1.2): 제목 쪽과 초록, 서론, 결과, 고찰, 방법(코드 가용성으로 끝남), 자료 가용성, 참고문헌,
%% 감사의 글, 저자 기여, 추가 정보, 그림 설명, 표 1, 그림 1-7.
\documentclass[10pt]{article}

% ---------- 레이아웃 ----------
\usepackage[a4paper,top=22mm,bottom=22mm,left=22mm,right=22mm]{geometry}
\usepackage{amsmath}
\usepackage{amssymb}
\usepackage{fontspec}
\usepackage{xeCJK}
\xeCJKsetup{CJKspace=true}
\usepackage{graphicx}
\usepackage{booktabs}
\usepackage{array}
\usepackage{xcolor}
\usepackage{catchfile}
\usepackage{titlesec}
\usepackage{caption}
\usepackage[numbers,sort&compress]{natbib}
\usepackage{url}
\usepackage[hidelinks]{hyperref}

% ---------- 글꼴 (본문 = Noto Serif CJK KR, 제목 = Pretendard) ----------
\setmainfont{Noto Serif CJK KR}[
  UprightFont={Noto Serif CJK KR}, BoldFont={Noto Serif CJK KR Bold},
  AutoFakeSlant=0.18]
\setCJKmainfont{Noto Serif CJK KR}[
  UprightFont={Noto Serif CJK KR}, BoldFont={Noto Serif CJK KR Bold},
  AutoFakeSlant=0.18]
\setCJKsansfont{Noto Sans CJK KR}
\setCJKmonofont{Noto Sans CJK KR}
\newfontfamily\headfont{Pretendard SemiBold}[BoldFont={Pretendard Bold}]
\newCJKfontfamily\cjkhead{Pretendard SemiBold}[BoldFont={Pretendard Bold}]

\linespread{1.3}
\setlength{\parindent}{1em}
\setlength{\parskip}{0.2em}

% ---------- 절 서식 (번호 없음, Sci Rep 영문판과 같음) ----------
\setcounter{secnumdepth}{0}
\titleformat{\section}{\headfont\cjkhead\large\bfseries}{}{0em}{}
\titleformat{\subsection}{\headfont\cjkhead\normalsize\bfseries}{}{0em}{}
\titlespacing*{\section}{0pt}{2.0ex}{0.8ex}
\titlespacing*{\subsection}{0pt}{1.4ex}{0.5ex}

% ---------- 캡션과 이름 ----------
\renewcommand{\figurename}{그림}
\renewcommand{\tablename}{표}
\renewcommand{\refname}{참고문헌}
\renewcommand{\abstractname}{초록}
\captionsetup{font=small,labelfont=bf,labelsep=period,justification=justified,singlelinecheck=false}

% ---------- 참고문헌 (영문판과 같은 sn-nature.bst, 위첨자 번호) ----------
\setcitestyle{super,open={},close={}}
\bibliographystyle{sn-nature}
\providecommand{\bibcommenthead}{}

% ---------- 그림과 표 1 (영문판과 같은 원본) ----------
\newcommand{\figdir}{../../../outputs/figures/paper/v3_restructure}
\CatchFileDef{\TableOneRaw}{\figdir/Table1_data.tex}{}
\long\def\ExtractTableOne#1\noindent\begin#2\bottomrule#3\EndExtractTableOne{%
  \def\TableOneTabular{\begin#2\bottomrule\end{tabular}}}
\expandafter\ExtractTableOne\TableOneRaw\EndExtractTableOne

\newcommand{\PlaceFigure}[2]{%
  \begin{figure}[p]
  \centering
  \IfFileExists{\figdir/#2}{%
    \makebox[\textwidth][c]{\includegraphics[width=166mm,height=0.88\textheight,keepaspectratio]{\figdir/#2}}%
  }{%
    \fbox{\parbox[c][80mm][c]{0.9\textwidth}{\centering\large 그림 #1 준비 중}}%
  }%
  \par\medskip
  {\small\bfseries 그림 #1}
  \end{figure}
  \clearpage}

\begin{document}

% 블록 1: 제목, 저자, 초록, 핵심어
\def\frontblock{1}% SI numbering: first-citation order (tools/si_numbering.py, round 6, 2026-10-05)
% front.tex : 국문판 제목 쪽, 초록, 핵심어(블록 1), 자료 가용성(블록 3), 감사의 글·저자 기여·추가 정보(블록 4).
% 영문판 paper/manuscript/en/sections/front.tex 의 같은 블록을 옮겼다(영문판 블록 2, Code availability 는
% 방법의 '코드 가용성' 소절을 쓴다). ko/main.tex 가 \frontblock 을 1, 3, 4 로 정한 뒤 블록마다 \input 한다.
% 제목·핵심어: MANUSCRIPT_SPEC 1.3 국문 열. 수치, 인용 키, 자리표시 문자열은 영문판과 같다.
\providecommand{\frontblock}{0}

% =====================================================================
% 블록 1. 제목 쪽, 초록, 핵심어
% =====================================================================
\ifnum\frontblock=1\relax
\title{희소 관측 영구동토 지역의 활동층 두께 예측을 위한 물리 기반 기계학습의 라벨 수별 워크플로}
\author{[DECISION: 저자, 소속, 교신 저자]\\{\small 교신 저자 전자우편: [DECISION: 교신 저자 전자우편]}}
\date{}
\maketitle

\begin{abstract}
\noindent 관측이 드문 영구동토 지역에서 Stefan 식과 기계학습(ML)으로 활동층 두께 지도를 만드는 라벨 수별 워크플로를 제시한다. ML 모델은 대개 학습 지역 안에서, 라벨 수 하나로, 재보정 기준선 없이 시험되어 왔다. 이 연구는 다섯 지역을 100~km 완충 구역과 함께 학습에서 빼고, 라벨 0개부터 전량까지 평균 제곱근 오차(RMSE)를 원천 계수와 재보정 계수의 Stefan 식과 비교하였다(내부 사전 등록). 라벨이 없을 때 Stefan 유사라벨은 섞은 유사라벨보다 오차가 낮았다($-1.6$~cm, 95\% 신뢰구간 [$-2.0$, $-1.0$]). 라벨 10개에서는 재보정이 원천 계수보다 4지역 평균 오차를 낮추었고($-2.5$~cm, [$-3.0$, $-1.9$], 오차가 낮아진 지역 1곳), 잔차 ML의 추가 차이는 확인하지 못했다. 이 결과를 본 뒤 설계한 분석에서 대상 라벨 안 교차검증으로 방법을 고르면 고정 잔차 레시피에 대한 RMSE 변화가 라벨 40개와 160개에서 $-0.87$~cm와 $-1.3$~cm였다(2지역, 변화는 캐나다에서 옴). 알래스카 안에서 잔차 ML의 RMSE는 13.9~cm, 재보정 Stefan 식은 14.4~cm였고(라벨 1000개), 오차 제곱 감소의 약 99\%(라벨 전량)는 기후 격자 셀 사이에 있었다. 라벨 100개를 블록에 흩어 놓을 때 무작위 추출에 대한 RMSE 변화는 캐나다 $-2.5$~cm, 레나델타 $+1.7$~cm였다. 다음으로 필요한 것은 격자 안 공변량이다. \textcolor{blue}{[PENDING: XE-a, XE-b, satellite inputs (stage 2); registered deadline 8 October 2026, 14:22 KST]}
% [DECISION: abstract XC sentence] 초록에서 XC 는 뺐다(FINAL_BATCH 7.1 은 수치만 쓰거나 빼는 것을 허용한다. 수치 문장은 S7 을 대신해도 영문 214단어가 된다).
% 후보 문장(영문판 후보와 같은 내용): 같은 지역을 다시 무작위로 나눈 분할에서 워크플로의 RMSE 변화는 라벨 160개에서 재보정 Stefan 식 대비
% −1.51 cm였다(지역 풀 평균; 3지역 가운데 2지역은 0.5 cm 넘게 오차가 컸다).
\end{abstract}

% 핵심어 6개(MANUSCRIPT_SPEC 1.3 국문 열). [DECISION: 핵심어]
\noindent\textbf{핵심어}: 활동층 두께, 영구동토, Stefan 식, 물리 기반 기계학습, 공간 전이, 라벨 배치
\fi

% =====================================================================
% 블록 3. 자료 가용성(본문 끝, 참고문헌 앞)
% =====================================================================
\ifnum\frontblock=3\relax
\section*{자료 가용성}

주 지역의 활동층 라벨은 GTN-P/CALM 자료\cite{Streletskiy2025}(PANGAEA, \url{https://doi.org/10.1594/PANGAEA.972777}), 레나강 삼각주의 ALLena 자료\cite{Veremeeva2025}(PANGAEA, \url{https://doi.org/10.1594/PANGAEA.973813}), ABoVE 토양 수분·활동층 두께 자료 2판\cite{Moore2025}(ORNL DAAC, \url{https://doi.org/10.3334/ORNLDAAC/2369})에서 받을 수 있다. 추가 지역의 라벨은 공개된 자료\cite{DuE2026,Grosse2007,Palmtag2022,Talucci2024,Pohl2026,Makarieva2017,Liang2023,Walker2009,Scheer2024,Martin2023,Boike2024,Hammar2025,Zastruzny2024,Sannel2020}에서, 라벨 정의 시험에만 쓴 지온 유도 라벨은 티베트 고원 자료\cite{Fu2025}에서 가져왔다. 식별자와 약관은 보충 표~S12에 있다. 재배포 약관이 확인되지 않은 자료의 라벨은 주 결과에 쓴 약관 확인판에서 뺐고 재배포하지 않으며, 보충 표~S12에 적은 원 제공처에서 받을 수 있다. 이런 라벨을 포함한 북대서양 전체 판은 보충 정보에만 보고한다.

공변량은 ERA5-Land 월평균 자료\cite{munozsabater2021_era5_land}(Copernicus Climate Data Store, \url{https://doi.org/10.24381/cds.68d2bb30}), SoilGrids 2.0\cite{Poggio2021} [미확인: 접근 시점의 SoilGrids 판 표지], Copernicus DEM GLO-30(Copernicus DEM, Global and European Digital Elevation Model, \url{https://doi.org/10.5270/ESA-c5d3d65}; AWS 공개 자료 버킷 copernicus-dem-30m 의 타일), ESA CCI Permafrost 활동층 두께 제품 4.0판\cite{Westermann2024}(CEDA, \url{https://doi.org/10.5285/d34330ce3f604e368c06d76de1987ce5})에서 만들었다. 알래스카 지역 내 시험의 합성 개구 레이더(SAR) 공변량은 2017년 ABoVE 항공 L·P 밴드 SAR 활동층 추정(ORNL DAAC, \url{https://doi.org/10.3334/ORNLDAAC/2004})과 알래스카 북부의 P 밴드 편파 SAR 활동층 상향 지도\cite{Whitcomb2024}에서 가져왔다 [미확인: PolSAR 상향 ALT 자료의 DOI, 후보 10.3334/ORNLDAAC/2332]. 지도 가림에는 ESA CCI Permafrost 범위 제품 4.0판(영구동토 비율, 1997--2021년 연별 파일) [미확인: 영구동토 범위 자료 v4.0 의 DOI]과 JRC Global Surface Water 출현 빈도 층\cite{Pekel2016}을 썼다.

조립한 라벨·공변량 표, 실행 산출물, 집계·판정 표, 셀 단위 예측, 모든 그림의 원자료는 [DECISION: 저장소와 DOI]에 올릴 예정이다. 약관이 확인되지 않은 자료의 행은 올리는 표에서 뺀다. [DECISION: KPDC 자료, SI 포함 여부와 식별자, 이용 조건]
\fi

% =====================================================================
% 블록 4. 참고문헌 뒤
% =====================================================================
\ifnum\frontblock=4\relax
\subsection*{감사의 글}

[DECISION: Acknowledgements, Author contributions, Competing interests 문구] 자료 가용성에 적은 자료 제공처에 감사한다. 활동층 두께 공변량은 ESA Climate Change Initiative Permafrost\_cci 과제의 자료를 썼다. 지형 공변량은 Copernicus WorldDEM-30 \textcopyright{} DLR e.V. 2010--2014와 \textcopyright{} Airbus Defence and Space GmbH 2014--2018로 만들었으며, 이 자료는 유럽연합과 ESA가 COPERNICUS로 제공하고 모든 권리는 보유자에게 있다. [미확인: ERA5-Land(Copernicus Climate Change Service) 사용 허가가 요구하는 귀속 문구] [DECISION: KPDC 자료, SI 포함 여부와 식별자, 이용 조건] [DECISION: 대회 보고서 인용 여부] [DECISION: 연구비 지원 기관과 과제 번호, 또는 지원 없음]

\subsection*{저자 기여}

[DECISION: Acknowledgements, Author contributions, Competing interests 문구] [DECISION: 저자, 소속, 교신 저자]

\subsection*{추가 정보}

이해 상충: [DECISION: Acknowledgements, Author contributions, Competing interests 문구]

교신과 자료 요청은 [DECISION: 저자, 소속, 교신 저자]에게 한다.
\fi


\section{서론}\label{sec:intro}
% intro.tex : 서론(국문판). 영문판 paper/manuscript/en/sections/intro.tex(2026-10-05 2차 수정판)를 문장 단위로 옮겼다.
% 수치와 인용 키는 영문판과 같다. 용어: MANUSCRIPT_SPEC 1.4 국문 정본, docs/GLOSSARY.md.

활동층 두께(active-layer thickness, ALT)는 영구동토 위 지표층이 계절에 따라 녹는 최대 깊이이며, 그 변화는 영구동토 위에 지은 기반 시설에 영향을 준다\cite{Karjalainen2019,Ran2022b}. ALT는 제한된 지점에서만 측정되며, 그 가운데 상당수가 35년에 걸친 Circumpolar Active Layer Monitoring(CALM) 자료에 있다\cite{Streletskiy2025}. 이 연구에서 평가한 다섯 지역의 ALT 측정값(라벨) 수는 러시아 동부 30개부터 알래스카 13,606개까지다(표~1). Stefan 해\cite{Stefan1891}는 ALT를 계수 $E$와 융해 도일(thawing degree-days, TDD) 제곱근의 곱으로 예측하며, $E$는 현지 융해 깊이 측정으로 보정한다\cite{Nelson1997}. 라벨이 있는 셀에서 $E$의 기하 평균은 레나델타 1.31\,cm\,($^{\circ}$C\,d)$^{-1/2}$, 러시아 서부 2.51, 티베트 고원 11.03이다(그림~1b). 새 영구동토 지역(대상 지역)에는 측정이 적으므로, 그곳 ALT 지도의 오차는 다른 지역(원천 자료)에 맞춘 모델에 달려 있다.

Stefan 식이나 Kudryavtsev 식의 출력은 기계학습(machine learning, ML)의 입력 특징으로 쓰였고\cite{Pilyugina2025,WangG2025}, $E$는 공변량으로 예측되었으며\cite{Ran2022b,ZhangC2024}, 물리 규칙으로 만든 라벨로 영점 장막(zero curtain) 검출기를 학습시킨 연구가 있다\cite{Gay2026}. 과정 모델 사전 학습과 잔차 모델링은 다른 환경 분야에서 쓰인다\cite{Read2019,Jia2021,Willard2022}. Gautam 등은 알래스카 CALM 68개 지점에서 무작위 70/30 분할 한 번으로 랜덤 포레스트와 Stefan 식을 비교하였다(시험 결정계수 0.24와 0.54)\cite{Gautam2025}. 대규모 ALT 지도는 거리 블록으로 검증되었다\cite{Aalto2018,Karjalainen2019,Ran2022a}. 군집된 공간 자료에서는 무작위 교차검증이 지도 정확도를 과대평가할 수 있고\cite{Roberts2017,Ploton2020,MeyerPebesma2022}, 편향 없는 지도 정확도를 얻으려면 공간 교차검증보다 설계 기반 추론을 쓰는 확률 표본이 필요하다는 주장이 있다\cite{Wadoux2021}. 위의 ALT 평가는 모두 라벨 조건 하나를 쓴다. 이웃 분야에서는 학습에서 뺀 지역이나 현지 관측이 적은 조건에서 모델을 평가하였고, 일부는 보정된 과정 모델과 비교하였다. 그 대상은 호수 수온\cite{Read2019,Jia2021,Willard2021}, 계측소 없는 유역\cite{Pool2019,Feng2023}, 지하 온도\cite{OMalley2026}, 공간 위험 예측\cite{Portes2026}이다. 우리가 아는 범위에서, 영구동토 ALT나 연평균 지중온도 연구 가운데 학습 자료에서 뺀 지역의 오차를 대상 라벨 수의 함수로, 같은 라벨로 재보정한 Stefan 식과 비교해 보고한 연구는 없다.

새 지역에는 대상 라벨이 쌓이는 동안 Stefan 식과 ML을 결합하는 규칙과 라벨을 더하는 절차가 필요하다. 이 연구는 세 질문을 다룬다. (1) 대상 라벨이 없을 때 물리 정보는 ML에 무엇을 더하며, 그 이득은 어디에서 오는가? (2) 대상 라벨이 10개에서 수천 개로 늘 때 ML은 같은 라벨로 재보정한 Stefan 식에 얼마나 더하며, 그 이득은 어느 공간 규모에서 생기는가? (3) 라벨이 쌓일 때 방법은 어떻게 고르고, 새 라벨은 어떤 절차로 배치해야 하는가?

각 대상 지역은 100\,km 완충 구역과 함께 원천 자료에서 빼고, 그 지역 0.5$^{\circ}$ 블록의 절반에서 라벨을 뽑아 나머지 절반에서 오차를 채점하였다. 라벨 수는 0, 3, 10, 40, 160, 320, 1000개와 쓸 수 있는 전량이다(그림~1c,d). 라벨이 많은 지역 안에서는 앞선 결과를 본 뒤 설계한 블록 홀드아웃 시험으로 재보정 모델에 대한 ML의 이득을 측정하였다. 대비는 블록 부트스트랩 95\% 신뢰구간(confidence interval, CI)과 $\pm$0.5\,cm 동등 한계로 판정하였고, 전이 분석은 내부 사전 등록하였다. 대상 라벨이 없을 때 Stefan 유사라벨(Stefan 예측을 학습 라벨로 쓴 것)은 섞은 유사라벨보다 주 4지역(레나델타, 캐나다, 러시아 서부, 러시아 동부)의 평균 제곱근 오차(root-mean-square error, RMSE) 평균을 낮추었다. 라벨 10개로 Stefan 계수를 재보정하면 원천 계수보다 이 값이 낮아졌다(오차가 낮아진 지역은 4곳 가운데 1곳). 알래스카 안에서 라벨 전량을 쓸 때 재보정 모델에 대한 ML 이득의 거의 전부는 기후 격자 셀 사이의 편향 보정에서 왔다. 이 연구의 기여는 평가 설계, 음성 결과를 포함한 그 결과, 그리고 그 결과로 만든 라벨 수별 워크플로다.


% SI numbering: first-citation order (tools/si_numbering.py, round 6, 2026-10-05)
% results_a.tex : 결과 R1-R4(국문판). 영문판 paper/manuscript/en/sections/results_a.tex(2026-10-05 2차 수정판)를
% 문장 단위로 옮겼다. 수치, 인용 키, 그림·보충 표 번호, 자리표시 문자열은 영문판과 같다.
% 신뢰구간은 [하한, 상한]으로 적는다(값은 영문판의 'a to b' 와 같다).

\section{결과}\label{sec:results}

\subsection{평가 설계와 검증 사다리}\label{sec:results-design}

앞에서 설명한 홀드아웃 절차(그림~1c,d)에서 모든 방법을 같은 대상 라벨을 쓰는 두 Stefan 기준선과 비교하였다(그림~2a,b). 하나는 원천 셀에서 계수 $E_0$를 추정한 원천 계수 Stefan 식(표~1)이고, 다른 하나는 대상 라벨로 계수를 다시 맞추되 $E_0$ 쪽으로 수축한 재보정 Stefan 식($\kappa = 10$)이다\cite{Steyerberg2004}.
4지역 평균은 레나델타, 캐나다, 러시아 서부, 러시아 동부에 같은 가중을 주며, 학습에서 뺀 참조 대상인 알래스카는 넣지 않는다.
구간은 블록 등가중이라고 적지 않으면 셀 가중이다.

다섯 지역을 같은 가중으로 평균하면 직접 ML의 RMSE는 무작위 셀 분할의 22.9~cm에서 지역 홀드아웃의 33.3~cm로 커졌다(차이 10.3~cm, 95\% CI [7.3, 12.8], 블록 등가중 6.1~cm [4.7, 7.6]). 같은 조건에서 학습 셀에 최소제곱으로 계수를 맞춘 Stefan 식은 26.8--27.0~cm에 머물렀다(그림~1e).
10.3~cm 증가는 확인적 결과가 아니라 효과 크기로 보고하며, 등록 대비인 500~km 완충 군집 홀드아웃과 무작위 셀의 차이는 8.2~cm였다(95\% CI [5.6, 10.4]; 보충 표~S1).
무작위 셀 분할에서 $k$겹 최근접 거리 정합(k-fold nearest-neighbour distance matching, kNNDM) 분할\cite{Linnenbrink2024}로 가면 직접 ML의 RMSE는 알래스카, 레나델타, 캐나다에서 5.76~cm, 8.17~cm, 14.20~cm 늘었고, 학습 셀에 맞춘 Stefan 식의 RMSE는 0.29~cm, 0.51~cm, 4.28~cm 늘었다(셀 가중; 서술, 앞선 결과를 본 뒤 설계; 그림~1e; 보충 표~S2).
확률 표본에서는 무작위 검증이 지도 정확도의 타당한 추정이라는 견해가 있으나\cite{Wadoux2021}, 이 연구의 라벨은 군집된 비확률 표본이므로 아래의 오차는 학습에서 뺀 블록에서의 예측 오차다.

\subsection{대상 라벨 없는 전이}\label{sec:results-label0}

대상 라벨이 없을 때 Stefan 식의 유사라벨은 섞은 유사라벨보다 RMSE를 1.6~cm 낮추었다(95\% CI [1.0, 2.0], 블록 등가중 [0.6, 1.5], 4지역, Holm 보정 \textit{P} = 0.001). 사후 집계에서 원천 계수 Stefan 식보다 오차가 2~cm 넘게 커진 대상은 30개 가운데 물리 유사라벨 증강에서 2개, 직접 ML에서 18개였다(그림~3b; 그림~4a,b).
두 상수 위약에 대한 감소는 1.7~cm와 3.6~cm였고, $\sqrt{\mathrm{TDD}}$에 선형인 위약에 대해서는 0.48~cm였다(등록 판정: 혼재).
2~cm 문턱은 전이 결과를 본 뒤 정하였고, 30개 대상에는 하위 지역 또는 라벨 확충판 23개가 들어 있으며, 구간에 근거한 집계도 같은 순서를 보였다(사후; 보충 표~S3).

주 구간과 공통 재표집 구간의 판정이 일치한 학습기 5종에서 직접 ML은 원천 계수 Stefan 식보다 4지역 평균 오차가 높았다(그림~4c).
주 학습기인 CatBoost에서 차이는 2.2~cm(95\% CI [1.1, 3.4], Holm 보정 \textit{P} = 0.023)였고, CatBoost의 다른 설정 2종과 TabPFN v2, TabICL v2에서는 0.85--1.9~cm였다.
이 다섯 판정은 분할 독립 가정에 의존하며, TabPFN v2의 판정은 보정 전에만 유의하였다(Holm 보정 \textit{P} = 0.117).
원천 계수 앵커에 낮은 가중($\lambda = 0.25$)을 준 잔차 모델은 학습기 3종에서 그 Stefan 식과 ±0.5~cm 안에서 동등하였다(분할 독립 가정에 의존).

알래스카를 학습에서 뺐을 때 물리 유사라벨은 원천 계수 모델보다 오차가 1.8~cm 높았다(95\% CI [0.4, 2.6]).
원천 계수가 지역 값과 멀리 떨어진 티베트 고원(그림~1b)에서는 직접 ML의 오차가 60.0~cm 낮았다(비맹검이고 교차 환경인 비교에서 예상된 결과).
라벨 10개에서 직접 TabICL v2는 4지역 평균 오차가 2.9~cm 낮았으나(95\% CI [1.1, 3.6]; 분할 독립 가정 의존), 알래스카를 포함한 3지역 평균에서는 6.0~cm 높았다.

같은 블록에서 제품 학습 지점으로부터 5~km 또는 25~km 안의 블록을 가린 뒤(등록 이탈; 앞선 결과를 본 뒤 설계), 라벨 없는 기존 CCI v5 지도는 원천 계수 Stefan 식보다 오차가 23.79~cm 높았고(Holm 보정 \textit{P} = 0.0006; 레나델타, 캐나다, 러시아 동부에서 높음), 잔차 ML은 같은 라벨 10개와 전량으로 재보정한 CCI v5보다 오차가 8.37~cm와 7.51~cm 낮았다(CCI v5 대비는 비맹검 부분 포함). 대상 지역 안의 CALM 지점도 일부 학습에 쓴 Wei v2 지도(제품 결측 대체 47.4\%)에서는 차이를 확인하지 못했다(ALT 정의가 서로 다름; 보충 표~S2).

\subsection{라벨 3--10개의 계수 재보정}\label{sec:results-recal}

대상 라벨 10개로 Stefan 계수를 재보정하면 원천 계수보다 4지역 평균 RMSE가 2.5~cm 낮아졌다(95\% CI [1.9, 3.0], 블록 등가중 [1.9, 3.3], Holm 보정 \textit{P} = 0.001). 오차가 낮아진 지역은 4곳 가운데 1곳(러시아 서부, 12.0~cm)이었고, 캐나다에서는 오차가 높아졌으며(2.3~cm, 95\% CI [0.7, 2.9]), 레나델타와 러시아 동부에서는 차이를 확인하지 못했다(지역 중앙값 변화 −0.06~cm, 러시아 서부 제외 +0.72~cm; 그림~2c--f; 그림~3a,d).
재보정 앵커에 잔차 ML을 더하면 RMSE가 추가로 −0.18~cm 달라졌다(95\% CI [−0.53, −0.01], Holm 보정 \textit{P} = 0.136). 이 차이는 보정 전에만 유의하고, 크기가 0.5~cm 미만이며, 분할 독립 가정에 의존한다.

라벨 전량에서 라벨을 쓰는 보정(다음 절의 선정 규칙)의 오차 감소는 라벨 10개로 잰 원천 계수 편향과 순위상관을 보였다(Spearman $\rho = 0.38$, 95\% CI [0.17, 0.55], 28대상, 단측 \textit{P} = 0.0002; 라벨 격자 결과를 본 뒤 설계).
편향의 부호는 감소와 관계가 없었다($\rho = -0.01$; 그림~5a,b; 보충 표~S4).
사후 분석에서 이 편향의 크기와의 관계는 재보정 이득, 선정 규칙이 더한 몫, 재보정을 넘어선 ML 몫의 어느 것에서도 확인하지 못했다(마지막 몫은 대상 고정 보조 구간에서만 양이었고, 약한 쪽 범주로 적었다)(그림~5d; 보충 표~S2).

라벨 10개 이하에서 잔차 구조는 Stefan 출력을 입력으로 받는 ML(물리 입력 ML)보다 4지역 평균 오차가 낮았다. 차이는 라벨 0개에서 2.7~cm, 라벨 10개에서 3.7~cm(95\% CI [2.8, 5.2], Holm 보정 \textit{P} = 0.001; 라벨 10개에서 4지역 가운데 2지역의 오차가 낮았고, 지역 중앙값 변화 −2.2~cm, 러시아 서부 제외 −0.32~cm)였다(그림~3c).
재보정 Stefan 값을 물리 입력으로 주어도 잔차 구조의 오차는 2.5~cm 낮았으나 캐나다에서는 3.1~cm 높았고, 이 지역 판정은 분할 독립 가정에 의존한다(보충 표~S5).

잔차 ML에 물리 유사라벨을 더한 결과는 라벨 10개에서 ±0.5~cm 안의 동등이었고, 라벨 40개의 레나델타와 캐나다에서는 오차가 0.34~cm 높았다(0.5~cm 미만; 보충 표~S5).
물리식과 제품을 라벨 가중으로 적층한 결과는 라벨 10개에서 재보정이나 잔차 ML과의 차이를 확인하지 못했고, 이 라벨 수에서 적층은 추출의 36\%에서 재보정으로 대체되었다(앞선 결과를 본 뒤 설계, 비맹검 부분 포함; 보충 표~S2).

\subsection{라벨 수십--수백 개의 방법 선정}\label{sec:results-selection}

대상 라벨 안에서 5겹 블록 교차검증으로 방법을 고르면 고정 잔차 레시피($\lambda = 0.25$)보다 라벨 160개에서 RMSE가 1.3~cm 낮아졌고(95\% CI [0.7, 1.7], 레나델타와 캐나다, 공통 재표집에서도 같음), 라벨 40개와 160개에서 비열등이었다(한계 0.5~cm). 두 결과 모두 캐나다에서 나왔으며, 이 규칙은 라벨 격자 결과를 본 뒤 설계하였다(그림~5c; 그림~6d).
라벨 40개에서의 감소 0.87~cm는 분할 독립 가정에 의존한다.
레나델타에서는 차이를 확인하지 못했다.
주 4지역의 라벨 10개에서는 비열등을 확인하지 못했고, 레나델타에서 규칙의 오차가 0.70~cm 높았다. 라벨 전량에서의 비열등은 분할 독립 가정에 의존한다.
규칙은 재보정 Stefan 식보다 오차가 낮은 대상의 수를 늘리지 않았다(보충 표~S4).

캐나다에서는 원천 계수 Stefan 식과의 차이를 확인하지 못했다(보충 표~S4).
알래스카를 학습에서 뺐을 때 규칙은 라벨 10, 40, 160개에서 원천 계수 모델보다 오차가 높았다(서술적 곡선 대비; 보충 표~S4).

라벨 40개와 160개에서 레나델타와 캐나다의 잔차 ML은 물리 입력 ML보다 오차가 1.9--2.4~cm 높았다(2.4~cm 대비는 분할 독립 가정에 의존)(그림~3c).
이 차이는 캐나다에서 생겼고, 알래스카를 학습에서 뺐을 때는 잔차 ML의 오차가 낮았다.
이 라벨 수에서 잔차 ML과 더 강한 재보정 기준선(토양 보정 Stefan 식을 $\kappa = 10$으로 재보정한 것) 사이의 차이는 확인하지 못했다.

원천 지역 하나 제외 교차검증으로 신경망을 조정한 결과는 학습기마다 달랐고, 학습기 전반에 대한 문장을 쓸 등록 조건을 채운 학습기는 없었으며, 이 조정은 대상 라벨 안의 선택과 비교하지 않았다(보충 표~S4).
같은 지역을 다시 무작위로 나눈 분할에서(재사용 지역의 재검정, 앞선 결과를 본 뒤 설계, 비맹검 부분 포함) 순차 워크플로는 무작위 배치와 고정 잔차 레시피보다 라벨 160개에서 오차가 0.93~cm 낮았고(레나델타와 캐나다; 캐나다 −2.41~cm, 레나델타 +0.54~cm 미결정), 라벨 40개에서는 차이를 확인하지 못했다(최소 검출 효과 약 0.5~cm; 레나델타 +0.21~cm). 워크플로의 배치 단계가 무작위 추출로 바뀌었으므로 이 대비는 같은 라벨에서 방법 규칙과 고정 레시피를 비교한 것이다.
같은 분할에서 워크플로는 라벨 10개, 40개, 160개에서 재보정 Stefan 식에 비열등이었고(라벨 10개는 5지역 평균, 40개와 160개는 레나델타와 캐나다 평균) 라벨 40개와 160개에서는 오차가 1.35~cm와 1.51~cm 낮았으나, 이는 풀 평균에 대한 것이다. 러시아 동부의 라벨 10개(+1.34~cm), 선택 계열인 알래스카의 라벨 40개와 160개(+1.51~cm와 +2.34~cm), 레나델타의 라벨 160개(+0.53~cm)에서는 0.5~cm 넘게 나빴다(그림~6e; 보충 표~S2).
지역 안의 보조 팔에서 규칙은 라벨 200개에서 잔차 ML과 동등하였고, 라벨 500개와 1000개에서는 차이를 확인하지 못했다.
적층은 라벨 40개와 160개에서도 차이를 확인하지 못했다(레나델타와 캐나다; 앞선 결과를 본 뒤 설계, 비맹검 부분 포함; 보충 표~S2).

% results_b.tex : 결과 R5-R8(국문판). 영문판 paper/manuscript/en/sections/results_b.tex(2026-10-05 2차 수정판)를
% 문장 단위로 옮겼다. 수치, 인용 키, 그림·보충 표 번호, 자리표시 문자열은 영문판과 같다.

\subsection{라벨이 많은 지역 안의 이득}
\label{sec:results-label-rich}

알래스카 안에서 재보정 Stefan 앵커 위의 잔차 ML은 같은 라벨로 재보정한 Stefan 식보다 라벨 1000개에서 RMSE를 0.54~cm 낮추었다(95\% CI [0.32, 0.69], 블록 등가중 [0.40, 0.80], 14.4~cm에서 13.9~cm, 공통 재표집에서도 같음). 라벨 전량에서 오차 제곱 감소의 약 99\%는 ERA5-Land 기후 격자 셀 사이에 있었다\cite{munozsabater2021_era5_land}(그림~7a,c).
이 시험은 앞선 결과를 본 뒤 설계하고 실행 전에 등록하였다(방법).

라벨 500개와 전량의 감소(0.39~cm와 0.52~cm)는 분할 독립 가정에 의존하며, 원천 계수 Stefan 식과의 차이는 라벨 1000개와 전량에서 확인하지 못했다(보충 표~S7).
레나델타와 캐나다에서는 같은 대비의 차이를 시험한 어느 라벨 수에서도 확인하지 못했다(보충 표~S7).

오차를 기후 격자 셀 사이와 안으로 나누면(대용: 0.5° 블록과 ERA5-Land TDD; 방법; 앞선 결과를 본 뒤 설계, 실행 전 등록), 라벨 전량의 잔차 ML은 재보정 모델보다 격자 사이 RMSE를 1.0~cm 낮추었고(95\% CI [0.6, 1.3]), 격자 안 RMSE는 ±0.5~cm 안에서 동등하였다(−0.01~cm, 95\% CI [−0.03, 0.04]). 토양 온도 격자 정의에서도 같았다(그림~7c).
격자 안 변동은 재보정 Stefan 식 오차 제곱의 72\%를 차지하였고, 잔차 ML은 그 가운데 알래스카에서 1\% 미만, 세 지역의 어느 설정에서도 8\% 미만을 설명하였다.
이득을 격자 규모의 편향 보정으로 보는 것은 사후 서술이다(채점 블록별 오차 변화는 그림~7b).

등록된 격자 안 가설에서 잔차 ML의 격자 안 예측은 라벨 전량에서 층화 지역 평균의 오차를 키웠다(+0.09~cm, 95\% CI [0.02, 0.40]). 토양 온도 정의와 캐나다에서도 같았고, 모두 통계적으로 구별되나 크기는 0.5~cm 미만이었다(보충 표~S10).
전이에서는 격자 안 예측이 관측 변동과 맞는다는 근거가 없었다.
더 따뜻한 블록으로의 외삽에 대한 등록 시험은 판정할 수 없었다(캐나다의 채점 블록 3개; 보충 표~S10).
다음 세 분석은 앞선 결과를 본 뒤 설계하였다(보충 표~S13).
같은 시기의 따뜻한 블록을 쓴 공간 대용(외삽 폭은 $\sqrt{\mathrm{TDD}}$에서 0.97~($^{\circ}$C~d)$^{1/2}$)은 실제 외삽 시험이 아니며, 여기서 잔차 ML은 재보정 Stefan 식보다 라벨 100개와 전량에서 오차가 2.44~cm와 2.41~cm 낮았고(비열등), 직접 ML과 두 모델 사이의 외삽 손실 차이는 확인하지 못했다.
지형, 수체, 수목 피복 입력 10개를 더해도 잔차 ML의 격자 안 오차는 모든 라벨 수에서 ±0.5~cm 안에서 동등하였다(보조 대비, 비맹검 부분 포함). \textcolor{blue}{[PENDING: XE-a, XE-b, satellite inputs (stage 2); registered deadline 8 October 2026, 14:22 KST]}
적층 잔차는 라벨 500개와 1000개에서 재보정 앵커 잔차보다 오차가 0.55~cm와 0.498~cm 낮았고(3지역 가운데 2지역; 뒤의 것은 0.5~cm 미만; CCI v5를 후보로 더하면 유지되지 않음), 라벨 200개와 전량에서는 차이를 확인하지 못했으며, 알래스카에서 적층은 재보정 Stefan 식에 비열등이었다(비맹검 부분 포함).

\subsection{새 라벨의 위치}
\label{sec:results-placement}

캐나다에서 라벨을 0.5° 블록이나 공변량 공간에 흩어 놓으면 무작위 셀 추출보다 RMSE가 전이 라벨 40개에서 2.4~cm(95\% CI [0.5, 3.4], 블록 등가중 [0.8, 2.5], 공변량 최대 최소 거리 추출), 지역 안 라벨 50--100개에서 2.5--3.2~cm 낮아졌다. 레나델타에서는 블록 층화가 라벨 100개에서 RMSE를 1.7~cm 높였다(95\% CI [0.9, 2.4], 블록 등가중 [0.1, 1.6])(그림~5).
모든 대비는 같은 분할에서 잔차 ML($\lambda = 0.25$)을 썼고, 앞선 결과를 본 뒤 설계하고 실행 전에 등록하였다.

전이 라벨 10개에서 어느 전략도 다섯 지역 가운데 오차가 높아진 곳이 없었으나, 이는 비열등 결과가 아니다. 한 가중에서는 레나델타의 블록 층화(+1.9~cm)와 러시아 동부의 공변량 분산(+0.36~cm)이 오차 증가를 보였기 때문이다(그림~5f; 보충 표~S9).

라벨 100개에서 알래스카 최선 전략인 공변량 군집 층화는 레나델타에서 RMSE를 높였다(+0.33~cm, 0.5~cm 미만).
내부 사전 등록한 블록 층화 대 무작위 셀 추출의 지역 평균 대비는 라벨 10개와 40개에서 미결정이었다(2단 재표집; 비맹검 부분 포함; 전이 대비와 같은 라벨 집합; 보충 표~S9).
평가 전에 고정한 배치 절차(방법)도 순차 워크플로 안에서 시험하였다(그림~7e).
같은 지역을 다시 무작위로 나눈 분할에서 이 절차는 셀 무작위 추출로 바뀌었으므로 배치 단계는 기여한 것이 없고, 그 효과는 시험하지 않았다(방법). 이 절차는 알래스카 계열에서 골랐고, 평가 지역의 셀은 배치와 방법 규칙의 설계에 쓰였으며, 평가 분할은 같은 셀을 새로 나눈 것이다.
블록 배분과 공변량 덮임을 합친 배치는 라벨 10개와 40개에서 무작위 배치와의 차이를 확인하지 못했고, 그 네 설정은 블록 또는 공변량 분산 배치와의 비교 16개 가운데 8개에서 비열등(2개는 보정 전에만 유의), 2개에서 오차가 높았으며(0.5~cm 미만), 레나델타를 겨냥한 사후 설계인 설계 가중 계수는 알래스카와 캐나다에서 오차를 키웠다(탐색; 앞선 결과를 본 뒤 설계, 비맹검 부분 포함; 보충 표~S13).
후보는 이미 측정한 셀이므로 현장 후보로 이득이 옮겨진다는 보장이 없다.

\subsection{독립 지역}
\label{sec:results-independent}

얕은 활동층을 가진 약관 확인 새 지역인 러시아 중부(라벨 57개)에서 대상 라벨 없는 직접 ML은 원천 계수 Stefan 식보다 RMSE를 2.7~cm 높였다(95\% CI [2.1, 4.1], 블록 등가중 [3.2, 6.6], 공통 재표집에서도 같음). 이 지역을 포함한 얕은 레짐 5지역 풀에서 라벨 전량의 잔차 ML은 같은 모델보다 RMSE를 3.2~cm 낮추었다(95\% CI [2.2, 4.2], 블록 등가중 [1.5, 3.7])(표~1).
[MISSING: 5지역 풀 대비(라벨 전량 잔차 ML − 원천 계수 Stefan, 라벨 10개 잔차 ML − 재보정 Stefan)의 지역 중앙값, 오차 감소·증가 지역 수, 러시아 서부 제외 값. 새 지역 판정표(약관 확인분 판)의 지역 행에서 수치 생성 스크립트로 만든다]
이 결과에는 교차 환경 표지가 붙는다. 풀이 두 계산 플랫폼에서 맞춘 지역을 합치기 때문이다(방법).
4지역 풀과 5지역 풀의 어느 지역에서도 라벨 0--40개의 직접 ML은 원천 계수 모델보다 오차가 낮지 않았다.
라벨 10개에서 잔차 ML과 재보정 Stefan 식은 5지역 풀에서 ±0.5~cm 안에서 동등하였으므로(−0.13~cm), 이 라벨 수의 4지역 결과는 이 풀에서 유지되지 않았다.
러시아 중부만 보면 라벨 전량의 잔차 ML은 차이를 확인하지 못했다. 두 가중의 판정이 엇갈렸기 때문이다(보충 표~S11).

깊은 활동층을 가진 티베트 고원은 따로 보고하며 독립 지역의 근거로 쓰지 않는다(채점 셀이 원천 고도 범위 밖에 있음; 비맹검 대비).
그곳에서 라벨 10개의 잔차 ML은 수축 재보정보다 RMSE가 20.9~cm 낮았으나, 같은 라벨에 최소제곱으로 맞춘 Stefan 식보다 낫지 않았다(+19.1~cm; 보충 표~S11).
자료 약관이 일부 확인되지 않은 북대서양 지역은 보충 표~S11에만 있다.
\textcolor{blue}{[PENDING: XF, new public regions; registered data deadline 11 October 2026, 14:22 KST; no eligible new macro region found so far]}
\textcolor{blue}{[PENDING: XC-F3, external-family test of the workflow; not tested so far, because placement reduced to random sampling makes its ten-label contrast identical, and no eligible new region; registered sentence after the XF data deadline, 11 October 2026, 14:22 KST]}
티베트 고원에서 지온 유도 보조 행은 라벨 3개와 10개로 재보정한 모델의 전이 오차를 8.14--83.17~cm 줄였다(앞선 결과를 본 뒤 설계, 비맹검 부분 포함). 이 라벨은 직접 측정과 정의가 달라 이 결과를 일반화하지 않으며, 원천 쪽 설계는 시험하지 않았다(약관 근거 미기록; 보충 표~S13).

\subsection{예측 구간}
\label{sec:results-intervals}

대상 라벨이 없을 때 지역을 묶음으로 보정한 계층 conformal 구간\cite{Dunn2023}은 학습에서 뺀 관측의 평균 0.86을 포함하였다(95\% CI [0.80, 0.92], 명목 0.90, 4지역). 보정 지역이 4개이므로 유한 표본 보장은 없다(앞선 규약으로 내부 사전 등록; 보충 표~S12).
알래스카에서 conformalized quantile regression으로 보정한 분위 구간\cite{romano2019_conformalized_quantile_regression}은 학습에서 뺀 네 지역에서 0.32--0.73을 포함하였다(보충 표~S12).
계층 구간의 폭을 정하는 다섯 방식에서 포함률은 0.86--0.95, 평균 폭은 81.3--132.4~cm였다(그림~6d).
구간 점수\cite{GneitingRaftery2007}에서는 셀별 흐름 척도와 지역 안에서 폭을 섞은 위약 사이, 그리고 어느 척도와 상수 폭 사이에서도 차이를 확인하지 못했다.
라벨 10개에서 계층 예측 분포와 보정된 상수 폭 구간 사이의 구간 점수 차이는 확인하지 못했다(+9.5~cm, Holm 보정 \textit{P} = 0.277).
라벨 40개와 160개에서 잔차 ML 구간은 레나델타, 캐나다, 알래스카에서 0.84--0.90을 포함하였다(그림~6e).


% SI numbering: first-citation order (tools/si_numbering.py, round 6, 2026-10-05)
% discussion.tex : 고찰(국문판). 영문판 paper/manuscript/en/sections/discussion.tex(2026-10-05 2차 수정판)를
% 문단과 문장 단위로 옮겼다. 수치, 인용 키, 자리표시 문자열은 영문판과 같다.

\section{고찰}\label{sec:discussion}

결과는 새 영구동토 지역을 위한 라벨 수별 워크플로를 정한다(그림~6d). 대상 라벨이 없을 때 기준은 원천 계수(또는 연도 정합) Stefan 식이며, ML은 물리 유사라벨이나 저가중 잔차로만 들어간다. 라벨 3--10개에서는 계수를 재보정하고(오차가 낮아진 지역은 4곳 가운데 1곳) 저가중 잔차를 더한다. 라벨 40개부터는 대상 라벨 안 교차검증으로 방법을 고르며, 그 이득은 캐나다에서 생겼다. 이 풀에서 라벨 3--160개에 권하는 방법의 점 추정 오차는 원천 계수 모델보다 높았다 [DECISION: Fig 7d 경로의 y 기준과 풀, 그림 명세 D-13]. 라벨이 수천 개인 지역 안에서는 잔차 ML이 뒤따르며, 그 이득은 기후 격자 셀 사이에 있다(그림~6a--c). 배치는 사전 지정한 절차로만 들어간다. 라벨을 흩어 놓는 것이 캐나다에서는 오차를 낮추었고 레나델타에서는 낮추지 못했기 때문이다(그림~7).

이득은 Stefan 계수 $E$의 수준 보정과 기후 격자 규모의 편향 보정에서 왔다. $E$는 토양 열전도도, n-인자, 토양 수분과 얼음의 잠열을 한데 묶은 계수이며\cite{Riseborough2008}, 그래서 지역 Stefan 지도는 토지 피복별 계수를 썼다\cite{Nelson1997,ShiklomanovNelson2002}. 이 연구에서 $E$는 융해 지수의 편향과 라벨 정의도 흡수한다(보충 주석~1). 라벨 전량에서 원천 계수에 대한 잔차 ML의 4지역 평균 감소 가운데 85\%는 재보정 몫이었고, 이 평균은 러시아 서부가 지배하였다(14.0~cm 가운데 재보정 13.5~cm; 사후 산술). 계수가 원천 값에 가장 가까웠으나 라벨 블록과 채점 블록 사이에서 달랐던 캐나다에서는 라벨 10개 재보정이 오차를 키웠고 선정 규칙이 대개 잔차 가중 1.0을 골랐으며, 우리는 이를 라벨 위치의 이질성을 보정한 것으로 해석한다.

라벨 0개 결과는 랜덤 포레스트가 Stefan 식보다 일반화가 나빴던 알래스카 무작위 분할 비교와 방향이 같다\cite{Gautam2025}. 앞선 알래스카 지역 내 분석 [DECISION: 대회 보고서 인용 여부]은 56개 설정 가운데 최선인 물리 잔차에서 이 연구의 더 엄격한 시험보다 큰 Stefan 식 대비 이득을 보고하였다(보충 주석~2). kNNDM\cite{Linnenbrink2024}은 이 군집 표본(알래스카 13,606행이 343개 위치에 있음)에서 지도와 라벨 사이의 거리를 흉내 낼 뿐이다.

다른 설명의 가능성은 낮다. Stefan 유사라벨은 섞은 유사라벨과 상수 유사라벨보다 오차가 낮았고(그림~3b), 선형 융해 지수 위약과의 차이는 0.5~cm 미만이었다. 따라서 이득은 제곱근 형태보다 셀 수준의 융해 지수 정보를 반영한다. 재보정 이득은 러시아 서부와 동부의 단년 라벨에서도 유지되었고, 캐나다에서는 점 라벨일 때 1~km 위치 평균일 때보다 0.86~cm 작았다(중앙값, 라벨 10개). 이 결과가 워크플로의 라벨 단위를 정한다. 물리식과 CatBoost 결과는 계산 환경 사이에서 일치하였고(최대 차이 $3.6\times10^{-14}$~cm) 신경망은 최대 5.6~cm 달랐으므로, 워크플로는 결정적인 이 방법들만 쓴다.

워크플로는 측정 지점의 순위를 매기지 않고 기후 격자 셀 안의 변동을 예측하지 않는다. 1~km 지도(보충 그림~S1--S5)는 표시 해상도이고, 그 정보 해상도는 기후 입력(약 10~km)과 토양 입력(약 5~km)에 묶인다. 결론은 지표 투과 레이더(ground-penetrating radar) 라벨을 포함한 원천 자료에 한정되며, 훈련 범위보다 따뜻한 기후는 평가할 수 없었다. 원천 자료의 알래스카 비중(원천 셀의 78--94\%)은 라벨이 적을 때 가장 크게 작용한다. 원천 계수가 재보정 계수에서 라벨 10개의 가중을 갖기 때문이다. 확인적 근거는 재사용한 주 지역 4곳과 얕은 레짐의 새 지역 1곳에서 온다. 비맹검인 티베트 고원은 유일한 깊은 레짐 지역이라 따로 보고하므로, 결과를 이 지역들 너머로 일반화하지 않는다. 값은 시험 설계에 따라 달라지므로(그림~1e), 무작위 분할 연구와는 직접 비교할 수 없다. 격자 셀 안의 토양 수분, 유기층 두께, 적설 기간, 미지형은 공변량에 거의 없으므로, 알래스카에서 설명한 격자 안 비율(1\% 미만)은 물리적 한계가 아니라 이 입력의 한계를 나타낸다. 현장 토양 수분도 격자 안 잔차를 설명하지 못했다(진단; 보충 표~S2). 방법 선정, 배치, 라벨이 많은 지역의 시험, 편향 상관, 이득 분해는 앞선 결과를 본 뒤 설계하였으므로 내부 사전 등록한 대비보다 약한 근거다.

각 단계에는 조건이 있다. 라벨은 1~km 위치 평균이고, 재보정은 라벨 10개 대비에 근거하며, 대상 라벨 안 교차검증 선정은 라벨이 덮는 0.5$^{\circ}$ 블록 수에 좌우되고 라벨 10개의 레나델타에서 고정 레시피보다 오차를 키웠다. 워크플로 전체는 같은 지역을 다시 무작위로 나눈 분할에서만 시험하였고, 그 배치 단계는 무작위 추출로 바뀌었으며, 비열등은 풀 평균에서만 성립하고 러시아 동부, 알래스카, 레나델타에서는 성립하지 않았다. 다음에 필요한 자료는 개발에 쓰지 않은 독립 지역 \textcolor{blue}{[PENDING: XF, new public regions; registered data deadline 11 October 2026, 14:22 KST; no eligible new macro region found so far]}, 같은 블록에서 채점한 기존 ALT 지도\cite{Ran2022a,Wei2026}(이 연구에서는 등록 이탈로 비교; 보충 표~S2), 그리고 무엇보다 기후 격자 셀 안에서 변하는 공변량이다. 알래스카에서 재보정 Stefan 식 오차 제곱의 72\%가 격자 셀 안에 있었다. \textcolor{blue}{[PENDING: XE-a, XE-b, satellite inputs (stage 2); registered deadline 8 October 2026, 14:22 KST]}


% 방법(마지막 소절은 코드 가용성)
% SI numbering: first-citation order (tools/si_numbering.py, round 6, 2026-10-05)
% methods.tex : 방법(국문판). 영문판 paper/manuscript/en/sections/methods.tex(2026-10-05 2차 수정판, XA 계열 수 정정 포함)를
% 문장 단위로 옮겼다. 수식, 수치, 인용 키, 커밋 번호, 자리표시 문자열은 영문판과 같다. 코드 경로는 원문 그대로 둔다.

\section{방법}\label{sec:methods}

\subsection{연구 지역과 라벨}\label{sec:m-labels}

주 지역의 활동층 두께(ALT) 라벨은 GTN-P/CALM 기록\cite{Streletskiy2025}, 레나강 삼각주 지역의 ALLena 자료\cite{Veremeeva2025}, 알래스카와 캐나다의 ABoVE 자료 2판\cite{Moore2025}에서 가져왔다. 한 위치의 반복 측정이 독립 라벨로 들어가지 않도록 기록을 셀로 모았다. 셀은 소수 넷째 자리로 반올림한 좌표(ABoVE와 ALLena)이거나 CALM 지점 하나이며, 셀의 라벨은 그 관측의 평균이다(보충 방법~1). 셀 표는 17,572행이고, 그 가운데 17,467행이 기계식 탐침이나 지표 투과 레이더(ground-penetrating radar, GPR)로 잰 직접 라벨이다. 메타자료에 지중 온도나 융해 관(thaw tube)이 적힌 CALM 셀 68개는 라벨로 쓰지 않는다. 지온에서 유도한 셀 20개(알래스카 15, 캐나다 4, 티베트 고원 1)는 표를 고정한 뒤에 발견되어 표지를 단 채 직접 라벨 범주에 남아 있다. 라벨 연도는 1990--2024년이므로, 예측은 정적이며 2015--2020년 공변량 기후값을 기준으로 한다.

광역 지역은 분석 전에 지리적으로 정하였으며, 알래스카, 레나델타, 캐나다, 러시아 서부, 러시아 동부, 러시아 중부, 그린란드다 [DECISION: 지역 이름 표기] [DECISION: 러시아 중부의 영문 표기]. 알래스카와 주 지역 4곳의 라벨 행 수는 러시아 동부 30개부터 알래스카 13,606개까지다(표~1). 하위 지역 10곳(보충 표~S11)은 라벨 없이 0.5$^{\circ}$ 블록 중심을 $k$-평균으로 군집해 정하였고, 본 실행 전에 고정하였다.

라벨 수 $n$은 행의 수이며, 행의 뜻은 지역마다 다르다. 행은 러시아 서부와 동부에서 CALM 지점의 다년 평균이고, 레나델타에서는 대개 한 번 방문한 관측이며, 알래스카와 캐나다에서는 대부분 점 관측이고, 추가 지역에서는 약 1~km 범위의 평균이다. 1~km 위치당 행 수는 알래스카 39.7, 레나델타 15.1, 캐나다 8.7이며, 알래스카 13,606행은 74개 블록의 343개 위치에 있다.

추가 지역은 어떤 적합보다 먼저 정한 규칙에 따라 공개 자료에서 모았다(보충 방법~2). 라벨은 1990년 이후의 계절 말 융해 깊이를 직접 잰 값이어야 하고, 기존 라벨 지역에서 100~km 이상 떨어져야 하며, 약 1~km 셀로 모았다. 유효 분할의 채점 블록 합집합이 8개 이상인 지역만 신뢰구간 산출에 적격으로 보았다. 재배포 약관이 확인되지 않은 자료는 주 판정을 정하는 약관 확인판에서 뺐다(보충 표~S12). 이 판에서는 러시아 중부(57행, 13블록, 얕은 레짐)와 티베트 고원(132행, 34블록, 깊은 레짐)이 적격이다. 북대서양 묶음(41행)은 전체 판에서만 적격이며 보충 정보에 보고한다.

\subsection{공변량}\label{sec:m-covariates}

모든 학습 방법은 모든 지역에서 쓸 수 있는 공변량 25개를 쓴다. 지형 공변량 6개는 고도, 경사, 사면 방향의 사인과 코사인, $33 \times 33$ 화소 창의 지형 위치 지수와 거칠기이며, Copernicus DEM GLO-30에서 계산하였다(자료 가용성). 기후 공변량 8개는 연평균 기온, 융해 도일(TDD), 동결 도일, $\sqrt{\mathrm{TDD}}$, 최난월과 최한월 기온, 1층 토양 온도, 적설 수량이며, 2015--2020년 ERA5-Land 월평균에서 얻었다\cite{munozsabater2021_era5_land}. 토양 공변량 9개는 5--15~cm 깊이의 점토, 모래, 실트 함량, 용적 밀도, 조립 파편, pH와 세 깊이의 토양 유기탄소이며, 250~m SoilGrids 2.0 층\cite{Poggio2021}을 약 5~km로 재표본해 추출하였다. 모델 기반 공변량 2개는 ESA CCI Permafrost ALT 제품 4.0판의 1997--2021년 평균과 그 유효 표지다\cite{Westermann2024}. 이 제품은 같은 재분석 자료로 일부 구동한 모델 산출물이므로 관측으로 다루지 않는다.

TDD와 Stefan 해\cite{Stefan1891}는 다음과 같다.
\begin{equation}
\mathrm{TDD} = \sum_{m} \max(\bar{T}_m, 0)\, d_m, \qquad \mathrm{ALT} \approx E\, s, \qquad s = \sqrt{\mathrm{TDD}},
\label{eq:stefan}
\end{equation}
여기서 $\bar{T}_m$은 $m$월의 2~m 기온 월평균, $d_m$은 그 달의 일수이고, $E$의 단위는 cm\,($^{\circ}$C\,d)$^{-1/2}$이다. 셀은 라벨, CCI 값, 1층 토양 온도로 계산한 TDD가 모두 유효할 때만 채점하며, 모든 방법이 같은 채점 셀을 쓴다. 관측 횟수나 연도처럼 라벨에서 유도한 양은 입력에서 빼며, 단위 시험으로 이를 강제한다. 입력의 정보 해상도는 기후가 약 0.1$^{\circ}$, 토양이 약 5~km다.

\subsection{Stefan 기준선과 계수 재보정}\label{sec:m-stefan}

계수 $E$는 토양 열 특성, 함수량, 지표 조건을 흡수하며, 지역 적용에서는 현지 측정으로 보정한다\cite{Nelson1997}. 대상의 원천 집합 $S$는 대상 지역과 그 둘레 100~km 완충 구역 밖의 모든 라벨 셀이며, 원천 계수는 원점을 지나는 최소제곱 기울기다.
\begin{equation}
E_0 = \frac{\sum_{i \in S} s_i\, y_i}{\sum_{i \in S} s_i^2},
\label{eq:e0}
\end{equation}
여기서 $y_i$는 관측 ALT다. 원천 계수 Stefan 식은 $E_0 s$를 예측한다. 주 대상의 원천 집합에서 알래스카 셀은 78--94\%를 차지하므로, $E_0$는 이 대상들 사이에서 거의 같다(1.59--1.62~cm\,($^{\circ}$C\,d)$^{-1/2}$; 표~1). 같은 추정량을 대상 라벨 $n$개로 계산한 값을 $E_{\mathrm{ls}}$라 하면, 재보정 Stefan 식은 다음 $E_n$으로 $E_n s$를 예측한다.
\begin{equation}
E_n = \frac{n\, E_{\mathrm{ls}} + \kappa\, E_0}{n + \kappa}, \qquad \kappa = 10.
\label{eq:en}
\end{equation}
따라서 원천 계수는 라벨 10개의 가중을 가지며, 이는 라벨이 적을 때 재보정이 오차를 키울 수 있기 때문이다\cite{Steyerberg2004}. $\kappa$ 값은 이 과제의 앞선 실험에서 정하였으므로 맹검이 아니며, $\kappa = 3$과 $\kappa = 30$은 민감도 사례다. 현지 최소제곱 Stefan 식은 수축 없이 $E_{\mathrm{ls}}$를 쓴다. 기계학습(ML)의 이득은 같은 라벨로 맞춘 재보정 모델에 대해서만 말한다. 원천 계수에 대한 이득에는 재보정의 이득이 들어 있기 때문이다.

물리 모델과 ALT 제품을 비교한 등록 분석에서 더 강한 기준선 두 가지를 더하였다. 연도 정합 Stefan 식은 각 라벨의 관측 연도(2010--2024년 안, 그 이전 라벨은 2010--2014년)에 대해 평균한 $\sqrt{\mathrm{TDD}}$에 $E_0$를 적용하며, 라벨 0개와 전량에서 보고한다. 라벨 40개와 160개에서는 토양 보정 Stefan 식(보충 방법~3)의 출력을 식~(\ref{eq:e0})과 식~(\ref{eq:en})으로 재보정해 더 강한 재보정 기준선으로 보고한다.

\subsection{기계학습 방법과 결합 구조}\label{sec:m-ml}

전이 설계에서 학습 행은 원천 셀과 대상 라벨 $n$개이며 가중은 같다. 앵커를 쓰는 방법은 다음 꼴이다.
\begin{equation}
\hat{y}(x) = a(x) + \lambda\, g(x),
\label{eq:anchor}
\end{equation}
여기서 $x$는 셀의 공변량, $a$는 물리 앵커, $g$는 그 잔차로 학습한 학습기다. 직접 ML은 $\sqrt{\mathrm{TDD}}$와 CCI 값을 포함한 공변량으로 ALT를 회귀한다. 물리 유사라벨 증강은 대상의 라벨 절반에 있는 셀에 값이 $E_n s$인 행을 원천 행 하나당 10개씩 더하며, 추출된 셀은 관측값을 유지한다. 원천 앵커 잔차 결합은 $a = E_0 s$를 쓴다. 재보정 앵커 잔차 결합(잔차 ML)은 $a = E_n s$를 쓰며, 잔차 목표는 원천 행에서 $y - E_0 s$, 대상 행에서 $y - E_n s$다. 증강·앵커·잔차 결합은 라벨 절반의 셀에 값이 0인 잔차 행을 더한다. 물리 입력 ML은 Kudryavtsev형 모델\cite{Kudryavtsev1974}과 토양 보정 Stefan 식의 출력, 또는 앵커 값을 직접 ML의 입력에 더한다. Stefan·CCI 앵커는 CCI 값이 유효한 곳에서 $E_n s$와 그 값의 평균이다. 본 실험의 다른 방법 네 가지는 보충 방법~3에 적었다. 잔차 가중 $\lambda \in \{0.25, 0.5, 1.0\}$은 적합 뒤에 적용한다. 판정의 기준 가중은 본 실행 전에 정하였으며, 잔차 구조에서 0.25, 직접 구조와 물리 입력 구조에서 1.0이다. 가중 0.25의 원천 앵커 잔차 결합을 저가중 잔차라 부른다.

위약 네 가지로 증강의 이득이 셀 수준 Stefan 값에서 오는지 시험하였다. 위약은 유사라벨을 셀 사이에서 섞은 같은 값, 대상 수준 평균, 원천 ALT 평균, 원천에 맞춘 TDD의 선형 함수로 바꾼 것이다.

방법 비교의 학습기는 CatBoost\cite{Prokhorenkova2018}다(반복 200회, 학습률 0.05, 깊이 3, seed 2개). 대상 라벨 없는 직접 ML에서는 학습기 10종을 비교하였다. 이 설정, 더 큰 CatBoost, 원천 지역 하나 제외 교차검증으로 조정한 CatBoost, 랜덤 포레스트, 미세 조정 없는 TabPFN v2\cite{Hollmann2025}와 TabICL v2\cite{Qu2026}, 그리고 원천 자료에서 원천 지역 하나 제외 교차검증으로 초모수를 고른 신경망 4종이다\cite{BergstraBengio2012,CawleyTalbot2010}. 신경망은 다층 퍼셉트론, 다중 헤드 퍼셉트론(TabM\cite{Gorishniy2025}과 다름), 축소 FT-Transformer형 모델\cite{Gorishniy2021}, RealMLP\cite{Holzmuller2024}다(설정은 보충 방법~5). 그림에서는 이 방법들을 Source Stefan, Year-matched Stefan, Recalibrated Stefan, Anchor + residual ML, Source anchor + residual, Physics pseudo-labels, Direct ML, Physics-input ML, Stefan-CCI anchor로, 아래에 적은 규칙은 Target-label CV selection으로 표시한다.

\subsection{전이 시험 설계}\label{sec:m-transfer}

가까운 라벨 셀에서 정보가 새지 않도록, 대상의 원천은 대상 지역과 100~km 완충 구역을 뺀다. 하위 지역은 상위 지역을 원천에서 뺀 방식과, 라벨이 있는 지역 안의 새 구역을 흉내 내도록 상위 지역의 나머지를 남긴 방식으로 평가하였다. 대상 27개는 알래스카를 뺀 3판 지역 6곳, 알래스카, 두 방식의 하위 지역 10곳이다.

각 대상의 0.5$^{\circ}$ 블록을 무작위로 순서를 바꾼 뒤 셀이 적은 쪽 절반에 차례로 배정하여 절반 A와 B로 다섯 번 나누었으므로, 라벨 셀과 채점 셀은 블록 경계로 떨어진다\cite{Roberts2017}. 라벨은 A에서 뽑고 오차는 B의 채점 셀에서 쟀다. A의 크기보다 작은 $\{0, 3, 10, 40, 160, 320, 1000\}$의 각 $n$과 A의 전체 셀에 대해, 대상, 방식, 분할, $n$, 추출로 정한 seed로 다섯 번 무작위 추출하여 모든 방법이 같은 라벨을 쓰게 하였다. 러시아 서부와 동부는 A의 셀이 14--17개이므로, 라벨 40개와 160개의 풀 평균에는 레나델타와 캐나다만 들어간다.

주 지역 4곳(레나델타, 캐나다, 러시아 서부, 러시아 동부)이 확인적 풀이다. 이 지역들은 이 과제의 앞선 실험에 쓰였으므로, 그 판정은 재사용 지역의 재검정이다. 알래스카는 참조 대상으로서 따로 보고하며, 레나델타, 캐나다와 함께 보조 3지역 평균에 넣는다. 3판 표의 러시아 중부와 그린란드는 점 추정으로만 보고하고, 상위 지역과 셀을 공유하는 하위 지역은 판정 범주별로 센다. 추가 지역은 얕은 레짐 5지역 확장 풀과 티베트 고원을 포함한 6지역 확장 풀에 들어간다. 티베트 고원의 채점 셀은 모두 원천 셀 고도의 99.5 백분위수보다 높으므로, 티베트 대비는 독립 지역의 근거로 쓰지 않는다.

검증 사다리는 알래스카와 주 지역 4곳에서 직접 ML과 학습 폴드에 맞춘 Stefan 식을, 무작위 셀부터 지역 홀드아웃까지 설계상 분리가 커지는 여덟 방식으로 채점하였다(보충 방법~3). 앞선 알래스카 분석은 $k$겹 최근접 거리 정합(kNNDM)\cite{Linnenbrink2024}을 더하였는데, 이 방식은 시험과 학습 사이의 거리를 예측 격자와 라벨 사이의 거리에 맞춘다.

\subsection{지역 내 시험 설계와 이득 분해}\label{sec:m-within}

지역 내 시험은 라벨이 수백에서 수천 개인 곳에서 ML이 무엇을 더하는지 묻으며, 전이 결과를 본 뒤 설계하였다. 알래스카, 레나델타, 캐나다에서 라벨은 절반 A에서 뽑았고 유일한 학습 행이었다. 전이 설계의 원천 계수를 사전값으로 썼고, $\lambda$는 뽑은 라벨 안의 블록 교차검증으로 $\{0.25, 0.5, 1.0\}$에서 골랐으며, 증강은 라벨 하나당 유사라벨 10개를 썼다. 첫 시험(분할 5개, 라벨 1000개 이상, 비교 상대는 교차검증으로 고른 최선의 물리 보정식, 알래스카에서는 SAR 공변량 9개 추가)은 기각되었다. 두 번째 시험은 네 요소를 바꾸었다. 비교 상대(재보정 Stefan 식), 분할(25개, 레나델타는 24개 유효), 레시피(Stefan·CCI 앵커와 직접 ML 추가), 라벨 범위(200, 500, 1000개와 전량)다. 두 시험을 모두 보고한다.

이득 분해는 ML이 어느 공간 규모에서 오차를 줄이는지 묻는다. 채점 셀을 0.5$^{\circ}$ 블록과 기온 $\sqrt{\mathrm{TDD}}$로 묶어, 한 묶음 안에서 Stefan 예측과 기후 공변량 8개가 일정하게 하였다. 묶음 $g$의 셀 수를 $N_g$, 평균을 $\bar{y}_g$와 $\bar{\hat{y}}_g$라 하면 다음과 같다.
\begin{equation}
\mathrm{SSE} = \sum_g \sum_{i \in g} \left[(y_i - \bar{y}_g) - (\hat{y}_i - \bar{\hat{y}}_g)\right]^2 + \sum_g N_g \left(\bar{y}_g - \bar{\hat{y}}_g\right)^2,
\label{eq:decomp}
\end{equation}
여기서 두 항은 격자 안 성분과 격자 사이 성분이다. 격자 안 RMSE는 셀이 2개 이상인 묶음을 쓰며, Stefan 식에서는 관측의 묶음 내 표준편차와 같다. 성분의 RMSE 차이는 더해지지 않으므로, 감소의 비율은 SSE로 계산한다. 등록 문서는 묶음 기준으로 토양 $\sqrt{\mathrm{TDD}}$를 적었으나 실행은 Stefan 예측의 기온 $\sqrt{\mathrm{TDD}}$를 썼다. 그래서 토양 기준 묶음의 민감도 분석을 등록하고 실행하였다(두 묶음 사이 수정 Rand 지수는 알래스카 0.999, 캐나다 0.997, 레나델타 0.832).

\subsection{라벨 배치 절차}\label{sec:m-placement}

배치 전략은 능동 선정을 빼고는 라벨 값을 쓰지 않고 절반 A의 셀을 고른다. 셀 무작위 추출 외의 전략은 블록 층화, 공변량 군집 층화, 공변량 최대 최소 거리 추출, 원천 외삽 우선, 잔차 학습기의 예측 분산에 따른 능동 선정, $\sqrt{\mathrm{TDD}}$ 우선이다(보충 방법~8). 지역 안에서는 라벨 20--500개에서 비교하였고, 블록 층화와 최대 최소 거리 추출은 전이에서 라벨 10개와 40개로도 비교하였다. 이 시험들은 결과를 본 뒤 설계하였고, 본 실험의 블록 분산 대비는 결과 전에 등록하였다.

알고리즘 P는 블록 배분과 공변량 범위를 결합한 절차이며, 워크플로 분석에서 시험한 사전 지정 배치 절차다. 입력은 후보 셀의 표준화 공변량(라벨 값 없음), 블록 크기 $N_b$, 라벨 10, 40, 160개의 중첩 예산, 배분 지수 $\gamma$, seed다.
\begin{enumerate}
\item seed로 고른 블록이나 기존 라벨이 있는 블록에서 시작해, 블록의 공변량 중심 $\mu_b$를 가장 먼 점 우선 순회로 정렬한다.
\item 예산 $n_k$에서 블록 $b$에 $q_b = n_k N_b^{\gamma} / \sum_{b'} N_{b'}^{\gamma}$개의 라벨을 최대 나머지 방식으로 배분하되 $N_b$를 넘지 않게 한다.
\item 남은 부족분이 가장 큰 블록을 채운다. 그 블록에 라벨이 없으면 $\mu_b$에 가장 가까운 셀을, 있으면 이미 고른 모든 라벨에서 가장 먼 셀을 고른다.
\item 예산이 늘어도 앞서 고른 라벨은 유지한다.
\end{enumerate}
후보 설정은 $\gamma = 0$, 0.5, 1, 블록의 첫 셀을 가장 먼 점 우선으로 고르는 변형, 블록 층화, 최대 최소 거리 추출이다. 평가 지역이 선택에 영향을 주지 않도록 알래스카 대상만으로 기계적 규칙에 따라 후보 하나를 고른다. 이 규칙에서 모든 후보가 제외되어 Algorithm P는 셀 무작위 추출이며, Algorithm P 대 무작위 배치의 워크플로 대비는 시험하지 않았다. 따라서 워크플로의 배치 단계는 기여한 것이 없다. 라벨 10개에서는 워크플로와 그 무작위 배치 비교 대상이 같은 라벨과 모델을 쓰고, 라벨 40개와 160개에서 무작위 배치와 고정 잔차 레시피에 대한 비교는 같은 라벨에서 선정 규칙과 고정 레시피를 비교한 것이다. 등록한 보고 규칙에 따라 측정 위치의 순위를 매긴 지도는 만들지 않았다.

\subsection{방법 선정 규칙과 편향 진단}\label{sec:m-selection}

선정 규칙은 대상 라벨만으로 방법을 고른다.
\begin{enumerate}
\item 라벨이 10개보다 적으면 재보정 Stefan 식을 쓴다.
\item 라벨이 있는 블록을 최대 5개 폴드에 무작위로 배정한다(블록이 하나면 셀 폴드). 폴드가 2개보다 적으면 재보정 Stefan 식을 쓴다.
\item 각 폴드에서 $E_n$을 다시 계산하고 후보 7개를 맞춘다. 후보는 원천 계수, 재보정, 현지 최소제곱 Stefan 식, $\lambda = 0.25$와 1.0의 잔차 ML, 증강·앵커·잔차 결합, 물리 유사라벨이다.
\item 교차검증 RMSE가 가장 낮은 후보를 고른다(동률이면 작은 $\lambda$와 앞선 후보).
\item 고른 후보를 모든 라벨로 다시 맞춘다.
\end{enumerate}
이 규칙은 라벨 격자 결과를 본 뒤 등록하였고, 대상 30개(전이 대상 27개와 약관 확인 추가 지역 대상 3개)의 라벨 10, 40, 160개와 전량에 적용하였다. 비교 상대는 고정 잔차 레시피이며 비열등 한계는 0.5~cm다.

편향 진단값은 한 추출의 라벨 10개에서 원천 계수 Stefan 식이 보인 평균 오차의 절댓값을 추출과 분할에 걸쳐 평균한 값이다. 재보정도 쓰는 라벨만 쓰므로, 어떤 학습기를 맞추기 전에 알 수 있다. 이 진단값과 라벨 전량에서 선정 규칙이 원천 계수 모델에 대해 보인 이득의 Spearman 상관을 대상 28개에서 계산하였다. 이 이득에는 재보정이 들어 있으므로, 상관은 라벨을 쓰는 보정(재보정과 잔차)에 관한 것이다. 재보정을 넘어선 이득의 몫은 따로 사후 분석으로 살폈다(추가 등록 분석).

\subsection{통계 분석}\label{sec:m-stats}

절반 B의 블록 $b$에서 오차 제곱합을 $\mathrm{SSE}_b$, 채점 셀 수를 $N_b$라 할 때, 두 가지 평균 제곱근 오차(RMSE, cm)를 계산하였다.
\begin{equation}
\mathrm{RMSE}_{\mathrm{cell}} = \left(\frac{\sum_b \mathrm{SSE}_b}{\sum_b N_b}\right)^{1/2}, \qquad \mathrm{RMSE}_{\mathrm{block}} = \frac{1}{|B|}\sum_b \left(\frac{\mathrm{SSE}_b}{N_b}\right)^{1/2},
\label{eq:rmse}
\end{equation}
여기서 $|B|$는 채점 블록 수다. 셀 가중 값이 주 값이고, 블록 등가중 값은 표본이 조밀한 블록의 영향을 줄인다. 대비 $\Delta$는 첫 방법의 RMSE에서 둘째 방법의 RMSE를 뺀 값으로, 분할, 추출, seed로 짝지어 분할에 걸쳐 평균하므로 음수는 첫 방법의 오차가 낮다는 뜻이다. 층화 평균은 지역에 같은 가중을 준다.

신뢰구간(CI)은 각 분할 안에서 채점 블록을 부트스트랩한 95\% 백분위 구간이며, 두 방법에 같은 재표집 블록을 쓴다. 라벨 격자 하네스가 모은 가설(본 실험과 추가 지역 실행)과 예측 구간 실험은 반복 1000회를, 다른 모든 대비는 10,000회를 썼다. 두 팔이 서로 다른 라벨 집합을 쓰는 본 실험의 블록 분산 대비와 워크플로 분석은 추출도 재표집하는 2단 부트스트랩을 썼고, 배치 시험은 라벨 추출을 조건으로 한 구간을 보고한다. 보조 구간은 한 지역 블록의 합집합을 모든 분할에 대해 한 번 재표집하였다. 이 구간이 판정을 바꾸는 경우 결과를 분할 독립 가정에 의존한다고 적으며, 이 단서가 없는 문장은 두 판정이 일치해야 한다. 지역은 유효 분할의 채점 블록 합집합이 8개 이상일 때 풀 구간에 들어갔다.

대비는 셀 가중과 블록 등가중 CI의 상한이 모두 0보다 작으면 오차 감소, 하한이 모두 0보다 크면 오차 증가, 둘 다 아니고 네 한계가 모두 $\pm 0.5$~cm 안에 있으면 동등(유의수준 2.5\%의 두 단측 검정\cite{Schuirmann1987}), 그 밖에는 미결정('차이를 확인하지 못했다')이다. 한계는 어떤 결과보다 먼저 정하였다. 이 한계는 주 지역에서 원천 계수 Stefan 식의 셀 가중 RMSE(21.7--43.0~cm)의 약 1--2\%이며, 측정 오차가 아니라 관례다. 1.0~cm와 그 RMSE의 2\%를 한계로 한 경우는 민감도 사례다. $|\Delta| < 0.5$~cm인 오차 감소나 오차 증가는 크기가 0.5~cm 미만인 구별 가능한 차이로 보고한다. 비열등은 단측이며 두 상한이 모두 $+0.5$~cm보다 작아야 한다. 대비의 최소 라벨 수 $n^{*}$는 셀 가중 CI 상한이 0보다 작고, 분할의 3분의 2 이상과 추출의 75\% 이상에서 $\Delta < 0$인 가장 작은 $n$이다.

유의수준은 양측 $\alpha = 0.05$이며, 단측 0.025에 해당한다. $P$ 값은 두 가중의 양측 부트스트랩 $P$ 값 가운데 큰 값이며, 모든 등록 대비의 보정 전과 Holm 보정 $P$ 값은 보충 표~S1에 있다. 결과 전에 정한 대비 10개만 초록에서 판정어를 쓸 수 있다. 이들의 $P$ 값은 가족 크기를 10으로 고정한 Holm 절차\cite{Holm1979}로 보정하였고, 보정 $P$가 0.05보다 작을 때만 초록에 방향을 적는다. 그렇지 않으면 초록에는 차이를 확인하지 못했다고 적고, 본문에 '보정 전에만 유의'를 더한다. 지역 내 시험, 배치 시험, 선정 시험과 추가 분석의 결과는 초록에 수치로만 나온다. 다른 등록 가족에서 Holm 보정 $P$ 값은 보조 값이며, 워크플로 시험과 제품 비교에서만 문구를 정한다.

풀 대비는 4--7개 지역에 근거하므로 재표집 단위는 0.5$^{\circ}$ 블록(지역당 13--74개)이며, 모든 풀 평균에 지역 값과 범위를 함께 적는다. 지역 일반에 관한 문장을 위해 핵심 대비를 주 지역 4곳과 알래스카의 분할 50개로 다시 계산하고, 부호 검정, 부호 뒤집기 순열 검정, Hartung-Knapp 확률효과 평균\cite{HartungKnapp2001}으로 요약하였다. 이런 문장은 Hartung-Knapp 구간이 0을 벗어날 때만 쓰며, 지역이 4곳이면 가장 작은 양측 부호 검정 $P$는 0.125다. 순위상관은 Spearman 계수이며, 구간은 블록과 추출의 결합 재표집으로 얻었다. 오차 하한은 지형이 아닌 거친 공변량이 같은 셀 5개 이상의 묶음 안에서 관측의 합동 분산이며(알래스카 13,333셀), 효과 크기의 분모로만 쓴다.

앞선 결과를 본 뒤 설계한 분석은 그렇게 표시하고, 사후 서술에는 판정어를 쓰지 않는다. 등록 가설마다 맹검 표지 네 가지 가운데 하나를 단다. 맹검, 부분 비맹검(다른 규약에서 방향을 봄), 재현(비맹검)(같은 대비의 방향을 봄), 등록 시점에 결과 존재(열지 않음)다. 예측 구간은 경험적 포함률과, $a = 0.1$인 중앙 구간 $[L, U]$의 구간 점수 $(U - L) + (2/a)(L - y)_{+} + (2/a)(y - U)_{+}$로 평가하였다\cite{GneitingRaftery2007}. 계층 conformal 구간은 원천 지역을 묶음으로 썼으며\cite{Dunn2023}, 묶음이 2--4개이므로 유한 표본 보장이 없다(보충 방법~4). $\pm$ 뒤에 값을 적지 않는다.

\subsection{내부 사전 등록과 등록 이탈}\label{sec:m-registration}

가설, 대비, 판정 규칙, 허용 문장은 해당 결과를 가져오기 전에 과제의 git 저장소에 커밋하였다(내부 사전 등록). 커밋에는 제삼자 시간 인증이 없으며, 시각은 한국 표준시의 커밋 시각이다(전체 표는 보충 방법~6). 라벨 격자 실험(LG, L1--L8)은 2026년 9월 29일 15:45:46에 커밋 40be64c로 등록하였다. 그 확장(LGX), 파운데이션 모델과 추가 지역 실행(LGT, LGD), 구간 실험(LGU), TabICL과 조정 신경망 실행(LGF), 초록 묶음 AB1--AB10을 포함한 보고 규칙(WRAPUP, 316714c, 2026년 9월 30일 03:02:42)이 뒤따랐고, 모두 그날 07:20:22의 첫 결과 회수보다 앞섰다.

워크플로 실험은 이 결과를 본 뒤 설계하고 실행 전에 등록하였다. WF0--WF5는 4168331(2026년 10월 1일), WF6--WF8은 1b42b4d(2026년 10월 1일), WF9와 WF10은 8180632(2026년 10월 2일), 토양 기준 묶음 민감도는 7da0064(2026년 10월 2일)다. WF0은 사후 서술이고, WF5로 계획한 우선순위 지도는 만들지 않았다. 분석 XA--XJ는 어떤 적합보다 먼저 커밋 2678100(2026년 10월 4일 14:22:39)으로 등록하였다. 실행 전 등록 확인 시험으로 표시한 것은 XC-F3뿐이며, XA는 사후 분석이다. 각 분석을 본문에 둘지 보충 정보에 둘지는 결과 전에 정하였다.

등록 이탈과 사후 변경은 다음과 같다.
\begin{enumerate}
\item 이득 분해(WF9)는 등록한 토양 $\sqrt{\mathrm{TDD}}$가 아니라 기온 $\sqrt{\mathrm{TDD}}$로 셀을 묶었고, 토양 기준 묶음의 민감도 분석을 더하였다.
\item 첫 지역 내 시험(WF1-a)이 기각된 뒤 두 번째 시험(WF6)은 비교 상대, 분할, 레시피, 라벨 범위를 바꾸었다.
\item 제품 비교(XG)는 구간 가설(LGU-B2)에 조건을 둔 서술 표로 등록되었다. 그 가설이 미결정으로 끝난 뒤 XG는 Holm 보정 대비 6개로 다시 등록하였고 등록 이탈 표지를 단다.
\item 워크플로 분할(XC)은 seed 1--200이 이미 쓰였으므로(LGX-N4) seed 6--15에서 201--210으로 옮겼고, XC 판정어를 초록에 쓰자는 제안은 철회하였다.
\item AB4에는 결과 전 민감도 계산이 같은 양의 분할 분포를 보여 주었다는 주석을 단다.
\item 표를 고정한 뒤 발견한 지온 유도 셀 20개는 표지를 단 채 직접 라벨 범주에 남아 있다.
\item 기술적 수정 네 건(ff1be02, 2026년 9월 30일 11:44:30)은 첫 회수 뒤에, 그러나 어떤 결과도 풀거나 열기 전에 커밋하였으며, 가설, 대비, 판정 규칙을 바꾸지 않는다. [DECISION: 수용 여부, 방법 초안 183행]
\item CatBoost는 교차 환경 점검 (i)을 통과하지 못했으므로, 플랫폼을 섞은 확장 풀 행에는 교차 환경 표지를 단다.
\item 등록 규칙에 따라 측정 우선순위 지도는 만들지 않았다.
\item 기후 외삽 시험(WF10)은 캐나다의 외삽 집합이 블록 3개여서 판정할 수 없었다.
\item XA는 등록한 7개 계열이 아니라 8개 대상 계열로 재표집하였다. 28대상 집합의 북대서양 대상이 등록 계열 목록에 없었기 때문이다(봉인 결과와 함께 기록한 구현 이탈).
\end{enumerate}

\subsection{추가 등록 분석}\label{sec:m-xbatch}

분석 10개(XA--XJ)는 위의 규약을 따르며, 각 분석을 본문과 보충 정보 가운데 어디에 둘지는 결과 전에 정하였다(보충 방법~9; 결과는 보충 표~S2).

XA(본문, 사후)는 원천 계수 모델에 대한 이득을 재보정 몫, 남은 수축 몫, 재보정을 넘어선 ML 몫으로 나누고, 두 정의의 계수 오차와 각 몫의 순위상관을 추정하며, 구간은 대상 계열과 블록의 결합 재표집으로 얻는다. 등록 목록은 7개 계열이었고, 그 목록에 없는 북대서양 대상은 8번째 단일 대상 계열로 재표집하였으며, 등록한 7개 계열의 27대상은 결과 범주가 같은 등록 밖 민감도로 분석하였다. 결과는 0을 벗어난 구간이거나 '확인하지 못함'으로만 적는다. XB(보충)는 교차 적합으로 대상 라벨에서 물리 모델과 ALT 제품의 볼록 가중을 학습하며, 잔차 학습기가 있는 경우와 없는 경우를 본다. XC(본문)는 알고리즘 P, 라벨 10개 진단, 방법 규칙을 얕은 레짐 5지역의 새 분할에 차례로 적용한다. Holm 보정 가설 8개는 라벨 40개와 160개에서 워크플로를 무작위 배치와 고정 잔차 레시피에 비교하고, 라벨 10, 40, 160개에서 재보정 Stefan 식에 대한 비열등을 시험하며, 선정 규칙 아래에서 알고리즘 P를 무작위 배치와 비교하고, 새로 얻은 지역에서 라벨 10개의 워크플로를 시험한다. 평가 셀이 구성 요소 설계에 쓰였으므로, 마지막 가설을 뺀 모두가 부분 비맹검이다.

XD(보충)는 알래스카 대상에서 알고리즘 P의 설정을 고르고, 다른 계수 추정량을 시험하며, 학습된 배치 정책을 알래스카 밖에서 $P$ 값 없이 한 번 평가한다. XE(보충)는 지역 안 잔차 ML에 10--500~m의 공개 입력을 더하고, 격자 안 RMSE를 공변량 25개 모델과 재보정 Stefan 식에 비교한다. XF(보충, 맹검)는 새 공개 자료에 적격 규칙을 적용하고 적격 지역마다 핵심 대비를 되풀이한다. XG(본문)는 제품 학습 지점에서 5~km나 25~km 안의 블록을 가린 뒤 기존 ALT 지도를 같은 블록에서 채점하며, 원래 값은 원천 계수 모델에, 재보정 값은 잔차 ML에 비교한다. XH(본문)는 세 지역에서 무작위 셀, 지점, 블록, kNNDM, 지역 홀드아웃 검증에 따른 재보정 Stefan 식, 잔차 ML, 직접 ML의 오차를 판정어 없이 서술한다. XI는 약관 확인 캐나다 확충 표에서 기후 외삽 시험을 되풀이하고, XJ는 지온 유도 라벨을 저가중 보조 행으로 시험한다(둘 다 보충).

\subsection{계산 환경과 재현성}\label{sec:m-computing}

물리 모델과 CatBoost는 CPU에서 결정적으로 실행된다. 라벨 격자 실험과 그 확장은 CPU 64코어와 NVIDIA T4 GPU 4개를 가진 클라우드 노드(Rescale)에서, 워크플로 실험과 추가 분석은 한 프로세서 계열(AMD EPYC Milan 96코어, Milan-X 64코어)의 클라우드 CPU 노드에서 실행하였다. 파운데이션 모델, 조정 신경망, 구간 실험, 추가 지역 실행은 NVIDIA RTX 3090 GPU를 가진 공유 로컬 서버를 썼다. 패키지 판은 고정하였고(모든 플랫폼에서 CatBoost 1.2.10; 보충 표~S13), 각 산출물은 코드 해시, 설정 해시, 패키지 판, 장치를 기록한다.

물리식과 CatBoost 결과는 클라우드와 로컬 실행 사이에서 34,360개 키에 걸쳐 $3.6 \times 10^{-14}$~cm 안에서 일치하였고, 클라우드 노드 두 종류는 공통 키 11,814개에서 같은 오차를 냈다. 신경망은 플랫폼 사이에서 재현되지 않았다. FT-Transformer 키 288개 가운데 43개가 0.5~cm 넘게 달랐고(최대 5.61~cm), 한 플랫폼 안의 재적합은 같았다. 릿지 앵커 위의 CatBoost는 최대 0.46~cm 달랐다. 따라서 대비의 양쪽은 한 플랫폼에서 오고, 플랫폼을 넘는 풀 평균에는 교차 환경 표지를 달며, 워크플로는 CatBoost와 닫힌 형태의 물리식만 쓴다.

\subsection{지도와 색표}\label{sec:m-maps}

지도는 70$^{\circ}$~N에서 축척이 참인 북극 평사 투영을 쓰고 중심 경도는 각 설명문에 적으며, 티베트 고원은 Lambert 방위 정적 투영의 삽도로 넣었다. 육지는 Natural Earth(1:50m), 영구동토 구역은 ESA CCI Permafrost 범위 제품 4.0판(영구동토 비율, 1997--2021년 평균)에서 가져왔고, 지도는 Cartopy 0.23.0으로 그렸다. 과학 색표\cite{Crameri2018Zenodo,Crameri2020}를 썼다. ALT에는 oslo를 뒤집어 0.12--0.90 범위로, 오차 차이에는 0을 중심으로 한 broc을, 잔차에는 bam을, 구간 폭에는 acton을 뒤집어 0.05--0.95 범위로 썼다.

\subsection{대형 언어 모델 사용}\label{sec:m-llm}

[DECISION: LLM 사용 기재 문구]

\subsection{코드 가용성}\label{sec:m-code}

자료 준비, 실험, 통계, 그림, 시험에 쓴 코드는 [DECISION: 공개 저장소 주소와 공개 시점]에서 받을 수 있고, 판 DOI와 함께 보관한다 [DECISION: Zenodo DOI]. 주요 구성 요소는 라벨 격자 실험(\path{scripts/3_deep_learning/h40_label_grid.py}), 검증 사다리(\path{scripts/2_evaluation/h41_validation_ladder.py}), 확장 실험(\path{scripts/3_deep_learning/h42_label_grid_ext.py}), 워크플로 실험(\path{scripts/3_deep_learning/h54_workflow.py}), 초록 대비 집계기(\path{scripts/2_evaluation/h39_scenarios.py}), 사후 재분석(\path{scripts/2_evaluation/wf0_reanalysis.py}), 추가 분석(\path{x_*.py}), 단위 시험(\path{tests/}), 클라우드 작업 설정(\path{configs/rescale/})이며, 등록 문서와 그 개정 이력도 함께 둔다. 추가 분석은 catboost 1.2.10, scikit-learn 1.9.1, pandas 2.3.3, scipy 1.17.1, numpy 2.1.1을 고정한다. TabPFN v2와 TabICL v2의 사전 학습 가중치는 재배포하지 않으며 제공처에서 받는다.


% 블록 3: 자료 가용성(참고문헌 앞)
\def\frontblock{3}% SI numbering: first-citation order (tools/si_numbering.py, round 6, 2026-10-05)
% front.tex : 국문판 제목 쪽, 초록, 핵심어(블록 1), 자료 가용성(블록 3), 감사의 글·저자 기여·추가 정보(블록 4).
% 영문판 paper/manuscript/en/sections/front.tex 의 같은 블록을 옮겼다(영문판 블록 2, Code availability 는
% 방법의 '코드 가용성' 소절을 쓴다). ko/main.tex 가 \frontblock 을 1, 3, 4 로 정한 뒤 블록마다 \input 한다.
% 제목·핵심어: MANUSCRIPT_SPEC 1.3 국문 열. 수치, 인용 키, 자리표시 문자열은 영문판과 같다.
\providecommand{\frontblock}{0}

% =====================================================================
% 블록 1. 제목 쪽, 초록, 핵심어
% =====================================================================
\ifnum\frontblock=1\relax
\title{희소 관측 영구동토 지역의 활동층 두께 예측을 위한 물리 기반 기계학습의 라벨 수별 워크플로}
\author{[DECISION: 저자, 소속, 교신 저자]\\{\small 교신 저자 전자우편: [DECISION: 교신 저자 전자우편]}}
\date{}
\maketitle

\begin{abstract}
\noindent 관측이 드문 영구동토 지역에서 Stefan 식과 기계학습(ML)으로 활동층 두께 지도를 만드는 라벨 수별 워크플로를 제시한다. ML 모델은 대개 학습 지역 안에서, 라벨 수 하나로, 재보정 기준선 없이 시험되어 왔다. 이 연구는 다섯 지역을 100~km 완충 구역과 함께 학습에서 빼고, 라벨 0개부터 전량까지 평균 제곱근 오차(RMSE)를 원천 계수와 재보정 계수의 Stefan 식과 비교하였다(내부 사전 등록). 라벨이 없을 때 Stefan 유사라벨은 섞은 유사라벨보다 오차가 낮았다($-1.6$~cm, 95\% 신뢰구간 [$-2.0$, $-1.0$]). 라벨 10개에서는 재보정이 원천 계수보다 4지역 평균 오차를 낮추었고($-2.5$~cm, [$-3.0$, $-1.9$], 오차가 낮아진 지역 1곳), 잔차 ML의 추가 차이는 확인하지 못했다. 이 결과를 본 뒤 설계한 분석에서 대상 라벨 안 교차검증으로 방법을 고르면 고정 잔차 레시피에 대한 RMSE 변화가 라벨 40개와 160개에서 $-0.87$~cm와 $-1.3$~cm였다(2지역, 변화는 캐나다에서 옴). 알래스카 안에서 잔차 ML의 RMSE는 13.9~cm, 재보정 Stefan 식은 14.4~cm였고(라벨 1000개), 오차 제곱 감소의 약 99\%(라벨 전량)는 기후 격자 셀 사이에 있었다. 라벨 100개를 블록에 흩어 놓을 때 무작위 추출에 대한 RMSE 변화는 캐나다 $-2.5$~cm, 레나델타 $+1.7$~cm였다. 다음으로 필요한 것은 격자 안 공변량이다. \textcolor{blue}{[PENDING: XE-a, XE-b, satellite inputs (stage 2); registered deadline 8 October 2026, 14:22 KST]}
% [DECISION: abstract XC sentence] 초록에서 XC 는 뺐다(FINAL_BATCH 7.1 은 수치만 쓰거나 빼는 것을 허용한다. 수치 문장은 S7 을 대신해도 영문 214단어가 된다).
% 후보 문장(영문판 후보와 같은 내용): 같은 지역을 다시 무작위로 나눈 분할에서 워크플로의 RMSE 변화는 라벨 160개에서 재보정 Stefan 식 대비
% −1.51 cm였다(지역 풀 평균; 3지역 가운데 2지역은 0.5 cm 넘게 오차가 컸다).
\end{abstract}

% 핵심어 6개(MANUSCRIPT_SPEC 1.3 국문 열). [DECISION: 핵심어]
\noindent\textbf{핵심어}: 활동층 두께, 영구동토, Stefan 식, 물리 기반 기계학습, 공간 전이, 라벨 배치
\fi

% =====================================================================
% 블록 3. 자료 가용성(본문 끝, 참고문헌 앞)
% =====================================================================
\ifnum\frontblock=3\relax
\section*{자료 가용성}

주 지역의 활동층 라벨은 GTN-P/CALM 자료\cite{Streletskiy2025}(PANGAEA, \url{https://doi.org/10.1594/PANGAEA.972777}), 레나강 삼각주의 ALLena 자료\cite{Veremeeva2025}(PANGAEA, \url{https://doi.org/10.1594/PANGAEA.973813}), ABoVE 토양 수분·활동층 두께 자료 2판\cite{Moore2025}(ORNL DAAC, \url{https://doi.org/10.3334/ORNLDAAC/2369})에서 받을 수 있다. 추가 지역의 라벨은 공개된 자료\cite{DuE2026,Grosse2007,Palmtag2022,Talucci2024,Pohl2026,Makarieva2017,Liang2023,Walker2009,Scheer2024,Martin2023,Boike2024,Hammar2025,Zastruzny2024,Sannel2020}에서, 라벨 정의 시험에만 쓴 지온 유도 라벨은 티베트 고원 자료\cite{Fu2025}에서 가져왔다. 식별자와 약관은 보충 표~S12에 있다. 재배포 약관이 확인되지 않은 자료의 라벨은 주 결과에 쓴 약관 확인판에서 뺐고 재배포하지 않으며, 보충 표~S12에 적은 원 제공처에서 받을 수 있다. 이런 라벨을 포함한 북대서양 전체 판은 보충 정보에만 보고한다.

공변량은 ERA5-Land 월평균 자료\cite{munozsabater2021_era5_land}(Copernicus Climate Data Store, \url{https://doi.org/10.24381/cds.68d2bb30}), SoilGrids 2.0\cite{Poggio2021} [미확인: 접근 시점의 SoilGrids 판 표지], Copernicus DEM GLO-30(Copernicus DEM, Global and European Digital Elevation Model, \url{https://doi.org/10.5270/ESA-c5d3d65}; AWS 공개 자료 버킷 copernicus-dem-30m 의 타일), ESA CCI Permafrost 활동층 두께 제품 4.0판\cite{Westermann2024}(CEDA, \url{https://doi.org/10.5285/d34330ce3f604e368c06d76de1987ce5})에서 만들었다. 알래스카 지역 내 시험의 합성 개구 레이더(SAR) 공변량은 2017년 ABoVE 항공 L·P 밴드 SAR 활동층 추정(ORNL DAAC, \url{https://doi.org/10.3334/ORNLDAAC/2004})과 알래스카 북부의 P 밴드 편파 SAR 활동층 상향 지도\cite{Whitcomb2024}에서 가져왔다 [미확인: PolSAR 상향 ALT 자료의 DOI, 후보 10.3334/ORNLDAAC/2332]. 지도 가림에는 ESA CCI Permafrost 범위 제품 4.0판(영구동토 비율, 1997--2021년 연별 파일) [미확인: 영구동토 범위 자료 v4.0 의 DOI]과 JRC Global Surface Water 출현 빈도 층\cite{Pekel2016}을 썼다.

조립한 라벨·공변량 표, 실행 산출물, 집계·판정 표, 셀 단위 예측, 모든 그림의 원자료는 [DECISION: 저장소와 DOI]에 올릴 예정이다. 약관이 확인되지 않은 자료의 행은 올리는 표에서 뺀다. [DECISION: KPDC 자료, SI 포함 여부와 식별자, 이용 조건]
\fi

% =====================================================================
% 블록 4. 참고문헌 뒤
% =====================================================================
\ifnum\frontblock=4\relax
\subsection*{감사의 글}

[DECISION: Acknowledgements, Author contributions, Competing interests 문구] 자료 가용성에 적은 자료 제공처에 감사한다. 활동층 두께 공변량은 ESA Climate Change Initiative Permafrost\_cci 과제의 자료를 썼다. 지형 공변량은 Copernicus WorldDEM-30 \textcopyright{} DLR e.V. 2010--2014와 \textcopyright{} Airbus Defence and Space GmbH 2014--2018로 만들었으며, 이 자료는 유럽연합과 ESA가 COPERNICUS로 제공하고 모든 권리는 보유자에게 있다. [미확인: ERA5-Land(Copernicus Climate Change Service) 사용 허가가 요구하는 귀속 문구] [DECISION: KPDC 자료, SI 포함 여부와 식별자, 이용 조건] [DECISION: 대회 보고서 인용 여부] [DECISION: 연구비 지원 기관과 과제 번호, 또는 지원 없음]

\subsection*{저자 기여}

[DECISION: Acknowledgements, Author contributions, Competing interests 문구] [DECISION: 저자, 소속, 교신 저자]

\subsection*{추가 정보}

이해 상충: [DECISION: Acknowledgements, Author contributions, Competing interests 문구]

교신과 자료 요청은 [DECISION: 저자, 소속, 교신 저자]에게 한다.
\fi


\bibliography{references}

% 블록 4: 감사의 글, 저자 기여, 추가 정보
\def\frontblock{4}% SI numbering: first-citation order (tools/si_numbering.py, round 6, 2026-10-05)
% front.tex : 국문판 제목 쪽, 초록, 핵심어(블록 1), 자료 가용성(블록 3), 감사의 글·저자 기여·추가 정보(블록 4).
% 영문판 paper/manuscript/en/sections/front.tex 의 같은 블록을 옮겼다(영문판 블록 2, Code availability 는
% 방법의 '코드 가용성' 소절을 쓴다). ko/main.tex 가 \frontblock 을 1, 3, 4 로 정한 뒤 블록마다 \input 한다.
% 제목·핵심어: MANUSCRIPT_SPEC 1.3 국문 열. 수치, 인용 키, 자리표시 문자열은 영문판과 같다.
\providecommand{\frontblock}{0}

% =====================================================================
% 블록 1. 제목 쪽, 초록, 핵심어
% =====================================================================
\ifnum\frontblock=1\relax
\title{희소 관측 영구동토 지역의 활동층 두께 예측을 위한 물리 기반 기계학습의 라벨 수별 워크플로}
\author{[DECISION: 저자, 소속, 교신 저자]\\{\small 교신 저자 전자우편: [DECISION: 교신 저자 전자우편]}}
\date{}
\maketitle

\begin{abstract}
\noindent 관측이 드문 영구동토 지역에서 Stefan 식과 기계학습(ML)으로 활동층 두께 지도를 만드는 라벨 수별 워크플로를 제시한다. ML 모델은 대개 학습 지역 안에서, 라벨 수 하나로, 재보정 기준선 없이 시험되어 왔다. 이 연구는 다섯 지역을 100~km 완충 구역과 함께 학습에서 빼고, 라벨 0개부터 전량까지 평균 제곱근 오차(RMSE)를 원천 계수와 재보정 계수의 Stefan 식과 비교하였다(내부 사전 등록). 라벨이 없을 때 Stefan 유사라벨은 섞은 유사라벨보다 오차가 낮았다($-1.6$~cm, 95\% 신뢰구간 [$-2.0$, $-1.0$]). 라벨 10개에서는 재보정이 원천 계수보다 4지역 평균 오차를 낮추었고($-2.5$~cm, [$-3.0$, $-1.9$], 오차가 낮아진 지역 1곳), 잔차 ML의 추가 차이는 확인하지 못했다. 이 결과를 본 뒤 설계한 분석에서 대상 라벨 안 교차검증으로 방법을 고르면 고정 잔차 레시피에 대한 RMSE 변화가 라벨 40개와 160개에서 $-0.87$~cm와 $-1.3$~cm였다(2지역, 변화는 캐나다에서 옴). 알래스카 안에서 잔차 ML의 RMSE는 13.9~cm, 재보정 Stefan 식은 14.4~cm였고(라벨 1000개), 오차 제곱 감소의 약 99\%(라벨 전량)는 기후 격자 셀 사이에 있었다. 라벨 100개를 블록에 흩어 놓을 때 무작위 추출에 대한 RMSE 변화는 캐나다 $-2.5$~cm, 레나델타 $+1.7$~cm였다. 다음으로 필요한 것은 격자 안 공변량이다. \textcolor{blue}{[PENDING: XE-a, XE-b, satellite inputs (stage 2); registered deadline 8 October 2026, 14:22 KST]}
% [DECISION: abstract XC sentence] 초록에서 XC 는 뺐다(FINAL_BATCH 7.1 은 수치만 쓰거나 빼는 것을 허용한다. 수치 문장은 S7 을 대신해도 영문 214단어가 된다).
% 후보 문장(영문판 후보와 같은 내용): 같은 지역을 다시 무작위로 나눈 분할에서 워크플로의 RMSE 변화는 라벨 160개에서 재보정 Stefan 식 대비
% −1.51 cm였다(지역 풀 평균; 3지역 가운데 2지역은 0.5 cm 넘게 오차가 컸다).
\end{abstract}

% 핵심어 6개(MANUSCRIPT_SPEC 1.3 국문 열). [DECISION: 핵심어]
\noindent\textbf{핵심어}: 활동층 두께, 영구동토, Stefan 식, 물리 기반 기계학습, 공간 전이, 라벨 배치
\fi

% =====================================================================
% 블록 3. 자료 가용성(본문 끝, 참고문헌 앞)
% =====================================================================
\ifnum\frontblock=3\relax
\section*{자료 가용성}

주 지역의 활동층 라벨은 GTN-P/CALM 자료\cite{Streletskiy2025}(PANGAEA, \url{https://doi.org/10.1594/PANGAEA.972777}), 레나강 삼각주의 ALLena 자료\cite{Veremeeva2025}(PANGAEA, \url{https://doi.org/10.1594/PANGAEA.973813}), ABoVE 토양 수분·활동층 두께 자료 2판\cite{Moore2025}(ORNL DAAC, \url{https://doi.org/10.3334/ORNLDAAC/2369})에서 받을 수 있다. 추가 지역의 라벨은 공개된 자료\cite{DuE2026,Grosse2007,Palmtag2022,Talucci2024,Pohl2026,Makarieva2017,Liang2023,Walker2009,Scheer2024,Martin2023,Boike2024,Hammar2025,Zastruzny2024,Sannel2020}에서, 라벨 정의 시험에만 쓴 지온 유도 라벨은 티베트 고원 자료\cite{Fu2025}에서 가져왔다. 식별자와 약관은 보충 표~S12에 있다. 재배포 약관이 확인되지 않은 자료의 라벨은 주 결과에 쓴 약관 확인판에서 뺐고 재배포하지 않으며, 보충 표~S12에 적은 원 제공처에서 받을 수 있다. 이런 라벨을 포함한 북대서양 전체 판은 보충 정보에만 보고한다.

공변량은 ERA5-Land 월평균 자료\cite{munozsabater2021_era5_land}(Copernicus Climate Data Store, \url{https://doi.org/10.24381/cds.68d2bb30}), SoilGrids 2.0\cite{Poggio2021} [미확인: 접근 시점의 SoilGrids 판 표지], Copernicus DEM GLO-30(Copernicus DEM, Global and European Digital Elevation Model, \url{https://doi.org/10.5270/ESA-c5d3d65}; AWS 공개 자료 버킷 copernicus-dem-30m 의 타일), ESA CCI Permafrost 활동층 두께 제품 4.0판\cite{Westermann2024}(CEDA, \url{https://doi.org/10.5285/d34330ce3f604e368c06d76de1987ce5})에서 만들었다. 알래스카 지역 내 시험의 합성 개구 레이더(SAR) 공변량은 2017년 ABoVE 항공 L·P 밴드 SAR 활동층 추정(ORNL DAAC, \url{https://doi.org/10.3334/ORNLDAAC/2004})과 알래스카 북부의 P 밴드 편파 SAR 활동층 상향 지도\cite{Whitcomb2024}에서 가져왔다 [미확인: PolSAR 상향 ALT 자료의 DOI, 후보 10.3334/ORNLDAAC/2332]. 지도 가림에는 ESA CCI Permafrost 범위 제품 4.0판(영구동토 비율, 1997--2021년 연별 파일) [미확인: 영구동토 범위 자료 v4.0 의 DOI]과 JRC Global Surface Water 출현 빈도 층\cite{Pekel2016}을 썼다.

조립한 라벨·공변량 표, 실행 산출물, 집계·판정 표, 셀 단위 예측, 모든 그림의 원자료는 [DECISION: 저장소와 DOI]에 올릴 예정이다. 약관이 확인되지 않은 자료의 행은 올리는 표에서 뺀다. [DECISION: KPDC 자료, SI 포함 여부와 식별자, 이용 조건]
\fi

% =====================================================================
% 블록 4. 참고문헌 뒤
% =====================================================================
\ifnum\frontblock=4\relax
\subsection*{감사의 글}

[DECISION: Acknowledgements, Author contributions, Competing interests 문구] 자료 가용성에 적은 자료 제공처에 감사한다. 활동층 두께 공변량은 ESA Climate Change Initiative Permafrost\_cci 과제의 자료를 썼다. 지형 공변량은 Copernicus WorldDEM-30 \textcopyright{} DLR e.V. 2010--2014와 \textcopyright{} Airbus Defence and Space GmbH 2014--2018로 만들었으며, 이 자료는 유럽연합과 ESA가 COPERNICUS로 제공하고 모든 권리는 보유자에게 있다. [미확인: ERA5-Land(Copernicus Climate Change Service) 사용 허가가 요구하는 귀속 문구] [DECISION: KPDC 자료, SI 포함 여부와 식별자, 이용 조건] [DECISION: 대회 보고서 인용 여부] [DECISION: 연구비 지원 기관과 과제 번호, 또는 지원 없음]

\subsection*{저자 기여}

[DECISION: Acknowledgements, Author contributions, Competing interests 문구] [DECISION: 저자, 소속, 교신 저자]

\subsection*{추가 정보}

이해 상충: [DECISION: Acknowledgements, Author contributions, Competing interests 문구]

교신과 자료 요청은 [DECISION: 저자, 소속, 교신 저자]에게 한다.
\fi


\clearpage
% legends.tex : 그림 설명(국문판). 영문판 paper/manuscript/en/sections/legends.tex(2026-10-05; Fig 6 은 렌더판
% Fig6_legend.md 와 맞춘 판)를 옮겼다. 그림 안 글자는 영문이므로 방법의 그림 표시 이름(Source Stefan,
% Recalibrated Stefan, Anchor + residual ML 등)은 원문 그대로 적는다. 수치, 인용 키, 자리표시 문자열은 영문판과 같다.

\section*{그림 설명}

\noindent 그림 1. 연구 지역, Stefan 계수, 평가 설계.
(a) 약관이 확인된 라벨 표의 라벨 셀을 0.5° 블록으로 모았다. 원의 면적은 블록 안 1~km 라벨 위치 수에 비례하고, New regions는 라벨 표를 고정한 뒤 더한 지역이다(러시아 중부 확충판; 티베트 고원은 삽도; 표~1).
음영은 연속 영구동토(영구동토 비율 $\geq$90\%)와 불연속 영구동토(50--90\%)다(ESA CCI Permafrost fraction v4.0, 1997--2021년 평균). 사각형은 c와 d의 범위다.
(b) 지역별 라벨 셀의 Stefan 계수 $E = \mathrm{ALT}/\sqrt{\mathrm{TDD}}$(로그 축). ALT는 활동층 두께, TDD는 ERA5-Land\cite{munozsabater2021_era5_land} 기온 융해 지수다(°C~d, 2015--2020년 평균). 지역마다 최대 400셀과 중앙값, 사분위 범위를 보이며, Source coefficient는 그 지역을 학습에서 뺐을 때 원천 풀에서 맞춘 계수다.
(c) 지역 홀드아웃. 레나델타의 라벨 셀과 그로부터 100~km 안의 모든 셀(Buffer)을 학습 자료에서 뺀다.
(d) 같은 지역을 라벨 블록과 채점 블록으로 나눈 분할(5개 가운데 하나)과 라벨 블록에서 뽑은 라벨 셀 40개. 오차는 채점 블록에서만 잰다.
지역 내 시험은 원천 자료 없이 같은 블록 분할을 쓰며, 전이 결과를 본 뒤 설계하였다.
(e) 무작위 셀부터 지역 홀드아웃까지의 검증에서 RMSE를 채점 셀에서 가장 가까운 학습 셀까지의 거리 중앙값에 대해 그렸다(선은 3지역 평균, 기호는 셀 가중 95\% 구간을 붙인 지역 값, 속이 빈 기호는 원천 계수 Stefan).
무작위 셀 분할에서 kNNDM 분할로 가면 직접 ML의 RMSE는 5.76~cm, 8.17~cm, 14.20~cm, 재보정 Stefan의 RMSE는 0.29~cm, 0.51~cm, 4.28~cm 늘었다(알래스카, 레나델타, 캐나다; 셀 가중; 서술, 앞선 결과를 본 뒤 설계).
약관이 확인되지 않은 셀(캐나다 38, 북대서양 19, 러시아 서부 1)은 그리지 않았다.
북극 평사 투영(중심 경도 127°~E), 축척 막대는 70°~N 기준이며, c와 d도 같은 투영이다. 티베트 고원 삽도는 Lambert 방위 정적 투영이고 축척 막대가 따로 있다. 해안선은 Natural Earth 50~m(c와 d는 10~m), Cartopy 0.23.

\medskip
\noindent 그림 2. 라벨 10개에서 Stefan 재보정은 4지역 평균 오차를 낮추었고, 잔차 ML의 추가 차이는 확인하지 못했다.
오차 변화는 학습에서 뺀 각 지역의 채점 블록에서 한 방법의 평균 제곱근 오차(RMSE)에서 원천 계수 Stefan 식(0선)의 RMSE를 뺀 값이며, 음수는 오차가 낮다는 뜻이다.
Recalibrated Stefan은 대상 라벨 $n$개로 계수 $E$를 다시 맞추되 원천 계수 쪽으로 수축하고($\kappa = 10$), Anchor + residual ML은 가중 0.25의 CatBoost 잔차를 더하며, Direct ML은 물리 앵커 없는 CatBoost이고, Year-matched Stefan은 라벨 연도의 융해 도일을 쓴다.
(a) 주 4지역(레나델타, 캐나다, 러시아 서부, 러시아 동부) 평균, 라벨 10개까지.
(b) 후보 라벨이 40개 이상인 레나델타와 캐나다의 평균. All은 모든 라벨이며, Year-matched Stefan은 점 추정이다.
(c--f) 지역별. 후보 풀보다 큰 라벨 수(러시아 서부·동부 14--17셀, 캐나다 325--406셀)는 가렸고, 패널 e는 세로 범위가 따로다(−12.5~cm까지).
선은 분할 5개의 셀 가중 평균이고, 띠는 재보정 기반 두 방법의 95\% 블록 부트스트랩 신뢰구간(CI)이다(패널 a, b의 구성 고정 평균은 재표집 10,000회, 패널 c--f는 1000회).
라벨 10개에서 재보정은 4지역 RMSE를 2.5~cm 낮추었고(95\% CI [1.9, 3.0]~cm, Holm 보정 $P = 0.001$), 오차가 낮아진 곳은 러시아 서부뿐이었으며(12.0~cm, 95\% CI [10.5, 13.6]~cm), 캐나다에서는 오차가 높아졌다(2.3~cm, 95\% CI [0.7, 2.9]~cm). 잔차를 더하면 RMSE가 −0.18~cm 달라졌다(95\% CI [−0.53, −0.01]~cm, 보정 전 $P = 0.034$, Holm 보정 $P = 0.136$).
그린 모든 대비의 판정은 보충 표~S2에 있다.
라벨 격자 분석은 내부 사전 등록하였고 한 계산 플랫폼에서 실행하였다.

\medskip
\noindent 그림 3. 물리 유사라벨은 섞은 유사라벨보다 오차를 낮추었고, 더 나은 결합 구조는 라벨 수에 따라 달랐다.
(a) 러시아 서부 라벨 셀 31개로 그린 개념도. 융해 도일 제곱근에 대한 활동층 두께와, 원천 계수 Stefan 직선(Source Stefan, 1.59~cm~(°C~d)$^{-1/2}$), 뽑은 라벨 10개로 재보정한 직선(Recalibrated Stefan, 2.11), 잔차 ML이 학습하는 그 라벨들의 편차(Residual)를 보인다. 두께는 아래로 갈수록 커진다.
(b) 물리 유사라벨로 학습한 직접 ML에서 위약 유사라벨로 학습한 같은 모델을 뺀 값, 4지역 평균, 라벨 0개와 10개.
(c) 잔차 구조에서 물리 입력 구조를 뺀 값을 라벨 수에 대해 그렸다. 입력 집합은 두 가지다(Physics inputs는 두 물리 모델의 출력, Recalibrated inputs는 원천 계수와 재보정 Stefan 예측). 범례는 등록 라벨 수(0과 10)와 보조 라벨 수를 구분하며, 라벨 40개와 160개의 값은 레나델타와 캐나다의 평균이다.
(d) Stefan 계수의 현지 최소제곱 재보정에서 수축 재보정을 뺀 값(서술; 양수는 수축이 유리), 라벨 10개까지는 4지역 평균, 그 위는 레나델타와 캐나다 평균.
구간은 95\% 블록 부트스트랩 신뢰구간(CI; 0.5° 채점 블록 재표집 10,000회)이며, 범례에 표시한 대로 셀 가중과 블록 등가중이다(패널 d는 셀 가중). 회색 띠는 ±0.5~cm다.
패널 b에서 물리 유사라벨은 섞은 유사라벨보다 라벨 0개의 오차를 1.6~cm 낮추었고(95\% CI [1.0, 2.0]~cm, Holm 보정 $P = 0.001$), 선형 융해 지수 유사라벨과의 차이는 라벨 0개에서 0.48~cm, 라벨 10개에서 ±0.5~cm 안이었다.
패널 c에서 라벨 10개까지 잔차 구조의 오차는 두 가중 모두에서 2.5--3.7~cm 낮았고(라벨 10개의 물리 입력에서 Holm 보정 $P = 0.001$), 라벨 40개와 160개에서는 물리 입력 구조의 오차가 1.9--2.4~cm 낮았으며 이 차이는 캐나다에서 왔다(대비 4개 가운데 1개는 분할 독립 가정에 의존).
위약과 결합 구조 분석은 내부 사전 등록하였다.

\medskip
% 2026-10-05 다시 그린 그림 4(XA 패널 d, 02:49:59)의 Fig4_legend.md 와 맞춘 영문판 설명문을 옮겼다.
\noindent 그림 4. 라벨을 쓰는 보정의 오차 감소는 라벨 10개로 잰 원천 계수 편향과 순위상관을 보였다.
(a) 재보정이 원천 계수 Stefan 식보다 오차를 낮춘(Recalibration) 시험 라벨 수의 최솟값 $n^{*}$와, 잔차 ML이 재보정 모델보다 오차를 낮춘(Residual ML) $n^{*}$(앞 격자 값부터의 구간). 판정은 등록한 세 조건을 따른다(방법). Not reached는 시험한 가장 큰 라벨 수까지 $n^{*}$가 없는 대상이다.
행은 b의 라벨 10개 편향 순서이며, 하위 지역은 독립이 아니다.
(b) 대상 라벨 안 교차검증으로 고른 방법의 오차 변화(라벨 전량, 원천 계수 Stefan 식 기준)를 처음 라벨 10개에서 그 모델이 보인 평균 편향의 절댓값에 대해 그렸다(Spearman $\rho = 0.38$, 95\% CI [0.17, 0.55], 28대상; 부호 있는 편향은 $\rho = -0.01$).
티베트 고원은 축 밖에 있으며(편향 227.7~cm; 오차 변화는 b에서 −159.6~cm, d에서 −135.7~cm와 −5.9~cm), 북대서양은 점 추정이다.
(c) 교차검증 선정에서 고정 잔차 레시피($\lambda = 0.25$)를 뺀 값.
선정은 레나델타와 캐나다 평균의 라벨 40개와 160개에서 비열등이었고(한계 0.5~cm), 라벨 160개에서 오차가 1.3~cm 낮았다. 라벨 40개의 0.87~cm 감소와 라벨 전량의 비열등은 분할 독립 가정에 의존하며, 라벨 10개에서는 비열등이 성립하지 않았다.
이득은 캐나다에서 왔고(2.1~cm와 2.7~cm), 레나델타의 변화는 +0.35~cm와 +0.07~cm였다.
(d) 라벨 전량에서 재보정의 오차 변화와, 더 나은 재보정 Stefan 식을 넘어선 잔차 ML의 오차 변화를 같은 편향에 대해 그렸다.
편향과 재보정 오차 감소의 관계($\rho = 0.37$, 95\% CI [−0.11, 0.93])와 재보정을 넘어선 감소의 관계($\rho = 0.28$, [−0.19, 0.80]; 대상 고정 구간 [0.18, 0.54])는 확인하지 못했다(사후 분석).
구간은 0.5° 채점 블록(d에서는 대상 계열도)을 10,000회 재표집한 95\% 신뢰구간(CI)이며, 범례에 표시한 대로 가중한다. 회색 띠는 ±0.5~cm다.
b와 c의 분석은 앞선 결과를 본 뒤 설계하였다.

\medskip
\noindent 그림 5. 라벨을 블록과 공변량 공간에 흩어 놓으면 캐나다에서는 무작위 셀보다 오차가 낮았고 레나델타에서는 그렇지 않았다.
(a--c) 캐나다 라벨 100개를 지도 위에 적은 세 배치 절차로 한 번 뽑은 결과(지역 내 시험, 분할 1; 방법). 작은 점은 라벨 블록의 후보 셀이고, 원의 면적은 지점마다 고른 라벨 수에 비례한다(40~km보다 가까운 블록은 합침).
지도는 절차를 보여 줄 뿐 지점의 순위를 매기지 않는다.
(d) Anchor + residual ML(잔차 가중 0.25)의 무작위 셀 대비 오차 변화를 캐나다, 레나델타, 알래스카 안의 라벨 수에 대해 그렸다.
(e) 같은 대비를 전이에서 라벨 10개와 40개로 본 것.
구간은 95\% 블록 부트스트랩 신뢰구간(CI; 0.5° 채점 블록 재표집 10,000회, 라벨 추출 조건부)이며, 범례에 표시한 대로 셀 가중과 블록 등가중이다. 회색 띠는 ±0.5~cm다.
패널 d에서 두 분산 전략은 캐나다의 라벨 50개와 100개에서 오차를 2.5--3.2~cm, 라벨 200개에서 1.3--1.4~cm 낮추었다.
레나델타에서 블록 층화의 오차 변화는 라벨 20--500개에서 +1.5~cm부터 +2.0~cm까지였고, 라벨 100개에서는 두 가중 모두 오차가 높았다(+1.7~cm; 분할 독립 가정에 의존).
알래스카의 라벨 50--200개에서는 두 가중의 판정이 엇갈렸다(블록 등가중 −1.1~cm부터 −1.7~cm까지, 셀 가중 +0.28~cm부터 +0.74~cm까지).
분산 기반 능동 선정은 알래스카와 레나델타의 라벨 50--200개에서 오차가 0.69--2.8~cm 높았고(대비 6개 가운데 1개는 분할 독립 가정에 의존), 캐나다에서는 흩어 놓기보다 이득이 작았다.
패널 e에서 공변량 분산은 캐나다의 라벨 40개에서 오차를 2.4~cm 낮추었다(95\% CI [0.5, 3.4]~cm; 분할 독립 가정에 의존). 라벨 10개에서는 두 가중 모두에서 오차가 높아진 지역이 없었으나, 레나델타 블록 층화의 셀 가중 구간은 0보다 위에 있었다(+1.9~cm).
배치 시험은 앞선 결과를 본 뒤 설계하고 실행 전에 등록하였다.
캐나다 중심의 평사 투영, 축척 막대는 패널 a. 해안선은 Natural Earth 50~m, Cartopy [미확인: 판].

\medskip
% 2026-10-05 05:3x: outputs/figures/paper/v3_restructure/Fig6_legend.md(04:59:09; Fig6.pdf 05:00:10, c 에 기존 ALT 지도 행)와 맞춘 영문판 설명문의 국문판.
\noindent 그림 6. 대상 라벨이 없을 때 2~cm 넘게 손해를 본 대상은 30개 가운데 물리 유사라벨에서 2개, 직접 ML에서 18개였다.
(a, b) 각 대상의 라벨 셀 중심에서 직접 ML(a)과 물리 유사라벨을 쓴 직접 ML(b)의 원천 계수 Stefan 식 대비 오차 변화(대상 17개; 원천에서 상위 지역을 뺌).
−40~cm보다 작은 값은 끝 색으로 칠한다(a의 티베트 고원, −60.0~cm).
알래스카 대상은 확대하였고, 레나델타 대상은 지시선으로 옮겨 그렸다.
제목의 개수는 대상 30개에 대한 사후 집계다(지역 안 하위 지역과 라벨 확충판 3개를 더함; 점 추정 문턱).
(c) 학습기별 직접 ML과 두 물리 기준선, 4지역 평균.
랜덤 포레스트, 물리 모델 입력을 받는 CatBoost, 원천 지역 하나 제외 교차검증으로 설정한 신경망 4종은 오차가 높았다(2.4--5.2~cm).
주 CatBoost 학습기(2.2~cm, 95\% CI [1.1, 3.4]~cm, Holm 보정 $P = 0.023$)를 포함한 나머지 학습기 5종도 주 구간에서 오차가 높았고, 이는 분할 독립 가정에 의존한다.
Year-matched Stefan은 오차가 1.0~cm 낮았고, 이것도 같은 가정에 의존한다.
Anchor ensemble은 척도 보정 Stefan, Kudryavtsev, ESA CCI 앵커의 평균이다.
기존 ALT 지도(등록 이탈; 앞선 결과를 본 뒤 설계; 제품 학습 지점 5~km 블록 제거): CCI v5는 오차가 높았고(+23.79~cm, 95\% CI [15.34, 30.16], 화살표), Wei v2(레나델타와 캐나다; 결측 대체 47.4\%)는 차이를 확인하지 못했다.
(d) 라벨이 없을 때 구간 척도 다섯 가지에서 90\% 예측 구간의 포함률을 평균 폭에 대해 그렸다. 4지역 평균(95\% 부트스트랩 구간, 재표집 1000회)과 지역 값이며, 겹쳐 적은 이름은 거의 같은 두 평균이고 위쪽이 포함률이 높다. 선은 명목 0.90이다.
(e) 라벨 40개와 160개에서 잔차 모델 구간의 포함률을 라벨 0개 대비 폭에 대해 그렸다. 띠는 0.85--0.95다.
c와 e의 구간은 95\% 블록 부트스트랩 신뢰구간(CI; 0.5° 채점 블록 재표집 10,000회)이며, c는 셀 가중과 블록 등가중이다(범례). ±0.5~cm 띠는 동등 한계다.
북극 평사 투영(중심 경도 127°~E; 알래스카 확대도 157°~W), 축척 막대는 70°~N 기준. 티베트 고원 삽도는 Lambert 방위 정적 투영.
영구동토 구역은 ESA CCI Permafrost fraction v4.0(1997--2021년 평균), 해안선은 Natural Earth 50~m, Cartopy 0.23.0.

\medskip
% 2026-10-05 05:3x: outputs/figures/paper/v3_restructure/Fig7_legend.md(05:27:56; Fig7.pdf 05:29:47, e = XC)와 맞춘 영문판 설명문의 국문판.
\noindent 그림 7. 알래스카 안에서 재보정 Stefan 식에 대한 잔차 ML의 이득은 1~cm 미만이었고, 기후 격자 규모의 편향 보정에서 왔다.
(a) 지역 내 시험(0.5° 블록의 절반을 채점, 분할 25개). 같은 라벨로 재보정한 Stefan 식 대비 Anchor + residual ML(교차검증 잔차 가중)과 Direct ML의 오차 변화.
Error floor는 공변량 조건부 하한(알래스카 11.31~cm, 레나델타 14.83~cm, 캐나다 17.90~cm)에서 재보정 Stefan RMSE를 뺀 값이며, 캐나다 선(−10.6~cm부터 −10.8~cm까지)은 축 아래에 있다.
알래스카의 라벨 1000개에서 잔차 모델은 오차가 0.54~cm 낮았다. 라벨 500개와 전량의 감소(0.39~cm와 0.52~cm)는 공통 재표집 구간에서 유지되지 않았다.
레나델타와 캐나다에서는 차이를 확인하지 못했다.
(b) 알래스카, 라벨 전량. 채점 블록별 오차 변화(분할 25개 합산)와 라벨 위치 343곳. 채점 셀이 10개 미만인 블록은 가렸다.
(c) ERA5-Land 기후 셀 사이와 안으로 나눈 오차 변화.
알래스카에서 오차 제곱 감소의 약 99\%는 셀 사이 성분에서 왔고(−1.03~cm; 셀 안 −0.01~cm), 3지역 평균에서는 셀 안 예측이 오차를 0.09~cm 키웠다.
(d) 라벨 구간별 권장 방법. 원천 계수 Stefan 식 대비 오차 변화, 레나델타와 캐나다 전이 풀(점 추정의 2지역 평균; 라벨 320개와 1000개는 시험하지 않음). 아래에는 배치와 진단 단계를 적었다. 라벨 40개부터의 선정은 고정 레시피에 대한 비열등에 근거한다(그림~4c).
(e) 같은 지역을 다시 무작위로 나눈 분할에서 워크플로(d; 배치는 무작위 추출로 바뀜)를 무작위 배치와 고정 레시피(−0.86~cm, −0.93~cm; 라벨 40개, 160개), 재보정 Stefan 식(+0.08~cm, −1.35~cm, −1.51~cm; 라벨 10개, 40개, 160개)과 비교하였다. 기호는 재보정 Stefan 식보다 0.5~cm 넘게 나쁜 지역이다. 앞선 결과를 본 뒤 설계, 재사용 지역의 재검정, 비맹검 부분 포함.
구간은 95\% 블록 부트스트랩 신뢰구간(CI; 0.5° 채점 블록 재표집 10,000회)이며, 셀 가중과 블록 등가중이다(범례). ±0.5~cm 띠는 동등 한계다.
패널 a--d는 전이 결과를 본 뒤 설계하였다.
알래스카 중심의 평사 투영, 축척 막대는 70°~N 기준. 해안선은 Natural Earth 50~m, Cartopy 0.23.0.

\section*{표 설명}

\noindent 표 1. 대상 지역, 라벨, 자료원.
Labels는 라벨 행 수이고, 1~km locations는 0.009° 셀 색인으로 행을 묶은 위치 수이며, 0.5° blocks는 라벨과 채점 분할의 단위다. Splits는 전이 설계의 유효 블록 분할 수이고, $E_0$는 그 지역을 학습에서 뺐을 때 원천 풀에서 맞춘 Stefan 계수(cm~(°C~d)$^{-1/2}$)이며, Alaska share of source cells는 원천 셀 가운데 알래스카에 있는 비율이다.
주 지역은 앞선 실험에 쓰였고 여기서 재검정하며, 알래스카는 주 지역과 평균하지 않는 참조 대상이고, 새 지역은 라벨 표를 고정한 뒤에 더하였다(약관 확인판).
하위 지역, 점 추정만 있는 대상, 북대서양 묶음은 보충 표~S1에 있다.
이용 조건은 자료 가용성에 적었다.


\clearpage
\begin{table}[p]
\caption{대상 지역, 라벨, 자료원. 표 본체는 영문판과 같은 생성 파일(Table1\_data.tex)이며 열 이름은 표 설명에 풀이하였다.}\label{tab:data}
\footnotesize\setlength{\tabcolsep}{4pt}%
\makebox[\textwidth][c]{\TableOneTabular}
\end{table}
\clearpage

\PlaceFigure{1}{Fig1.pdf}
\PlaceFigure{2}{Fig2.pdf}
\PlaceFigure{3}{Fig3.pdf}
\PlaceFigure{4}{Fig4.pdf}
\PlaceFigure{5}{Fig5.pdf}
\PlaceFigure{6}{Fig6.pdf}
\PlaceFigure{7}{Fig7.pdf}

\end{document}
